\begin{enumerate}
\def\labelenumi{(\alph{enumi})}
\setcounter{enumi}{1}
\tightlist
\item
  Suppose that G is compact. Let \(\phi\) be an étale endomorphism of G. Then \(\operatorname{Ker} \phi\) is finite, \(\phi(G)\) is of finite index in G and
\end{enumerate}

\[
\frac {\operatorname{Card} (\operatorname{Ker} \phi)}{\operatorname{Card} (\mathrm{G} / \phi (\mathrm{G}))} = \operatorname{mod} \det \mathrm{L} (\phi).
\]

(Use (a) to transform \(\mu (\mathbf{G})\) .)

\begin{enumerate}
\def\labelenumi{(\alph{enumi})}
\setcounter{enumi}{2}
\tightlist
\item
  Suppose that G is compact and commutative. Let \(r \in \mathbf{Z}\) be such that \(r, 1 \neq 0\) in K and let \(\phi\) be the endomorphism \(x \mapsto x^r\) of G. Then
\end{enumerate}

\[
\frac {\operatorname{Card} (\operatorname{Ker} \phi)}{\operatorname{Card} (\mathrm{G} / \phi (\mathrm{G}))} = (\mathrm{mod} r. 1) ^ {n}.
\]

(Observe that \(\mathrm{L}(\phi)\) is the homothety of ratio \(r\) by §2, Corollary to Proposition 7.)

\begin{enumerate}
\def\labelenumi{\arabic{enumi}.}
\setcounter{enumi}{2}
\tightlist
\item
  Let E be the complex vector space of symmetric complex matrices with 2 rows and 2 columns. For all \(s \in \mathbf{SL}(2, \mathbf{C})\) , we define an automorphism \(\varphi(s)\) of E by \(\varphi(s).m = s.m.^t s\) .
\end{enumerate}

\begin{enumerate}
\def\labelenumi{(\alph{enumi})}
\tightlist
\item
  Show that \(\rho\) is an analytic linear representation of \(\mathbf{SL}(2, \mathbf{C})\) . Identifying \(\begin{pmatrix} a & b \\ b & c \end{pmatrix} \in \mathrm{E}\) with \((a, b, c) \in \mathbf{C}^3\) , the matrix of \(\rho\begin{pmatrix} \alpha & \beta \\ \gamma & \delta \end{pmatrix}\) is
\end{enumerate}

\[
\left( \begin{array}{c c c} \alpha^ {2} & 2 \alpha \beta & \beta^ {2} \\ \alpha \gamma & \alpha \delta + \beta \gamma & \beta \delta \\ \gamma^ {2} & 2 \gamma \delta & \delta^ {2} \end{array} \right).
\]

\begin{enumerate}
\def\labelenumi{(\alph{enumi})}
\setcounter{enumi}{1}
\tightlist
\item
  Show that \(m \mapsto \det(m)\) is the non-degenerate quadratic form
\end{enumerate}

\[
q (a, b, c) = a c - b ^ {2}
\]

on E and that \(\rho\) is a morphism of \(\mathbf{SL}(2,\mathbf{C})\) onto \(\mathbf{SO}(q)\) with kernel \(\{I, - I\}\) . (For the surjectivity, observe that, in the matrix of an element of \(\mathbf{SO}(q)\) , the 1st column can be expressed in the form \((\alpha^2,\alpha \gamma ,\gamma^2)\) and the 3rd column in the form \((\beta^2,\beta \delta ,\delta^2)\) , where \(\alpha^2\delta^2 +\beta^2\gamma^2 -2\alpha \gamma \beta \delta = 1\) . Then observe that 2 elements of \(\mathbf{SO}(q)\) which coincide on a non-isotropic plane are equal.)

\begin{enumerate}
\def\labelenumi{(\alph{enumi})}
\setcounter{enumi}{2}
\tightlist
\item
  Deduce that the complex Lie group \(\mathbf{SO}(3, \mathbf{C})\) is isomorphic to the complex Lie group \(\mathbf{SL}(2, \mathbf{C}) / \{I, -I\}\) .
\end{enumerate}

\begin{enumerate}
\def\labelenumi{\arabic{enumi}.}
\setcounter{enumi}{3}
\tightlist
\item
  \begin{enumerate}
  \def\labelenumii{(\alph{enumii})}
  \tightlist
  \item
    Let \(q\) be a quadratic form of signature (1, 2) on \(\mathbf{R}^3\) . Let \(\mathbf{P}\) be the set of \(m \in \mathbf{R}^3\) such that \(q(m) > 0\) . Then \(\mathbf{P}\) is the union of two disjoint convex cones \(\mathbf{C}\) and \(-\mathbf{C}\) . If \(s \in \mathbf{SO}(q)\) , then either \(s(\mathbf{C}) = \mathbf{C}\) and \(s(-\mathbf{C}) = -\mathbf{C}\) or \(s(\mathbf{C}) = -\mathbf{C}\) and \(s(-\mathbf{C}) = \mathbf{C}\) . Let \(\mathbf{SO}^{+}(q)\) be the set of \(s \in \mathbf{SO}(q)\) such that \(s(\mathbf{C}) = \mathbf{C}\) . Then \(\mathbf{SO}^{+}(q)\) is an open normal subgroup of \(\mathbf{SO}(q)\) of index 2 in \(\mathbf{SO}(q)\) .
  \end{enumerate}
\end{enumerate}

\begin{enumerate}
\def\labelenumi{(\alph{enumi})}
\setcounter{enumi}{1}
\item
  By imitating the method of Exercise 3, define a morphism of \(\mathbf{SL}(2,\mathbf{R})\) onto \(\mathbf{SO}^{+}(q)\) with kernel \(\{I, - I\}\) . Deduce that the real Lie group \(\mathbf{SO}^{+}(q)\) is isomorphic to the real Lie group \(\mathbf{SL}(2,\mathbf{R}) / \{I, - I\}\) .
\item
  For \((\xi, \eta) \in \mathbf{C}^2\) and \((\xi', \eta') \in \mathbf{C}^2\) , let \(f((\xi, \eta), (\xi', \eta')) = \xi \bar{\xi}' - \eta \bar{\eta}'\) . Show that the inner automorphism of \(\mathbf{SL}(2, \mathbf{C})\) defined by \(\frac{1}{\sqrt{2}} \begin{pmatrix} 1 & -i \\ -i & 1 \end{pmatrix}\) maps \(\mathbf{SU}(f)\) to \(\mathbf{SL}(2, \mathbf{R})\) .
\end{enumerate}

\begin{enumerate}
\def\labelenumi{\arabic{enumi}.}
\setcounter{enumi}{4}
\tightlist
\item
  Let E be the real vector space of Hermitian complex matrices with 2 rows and 2 columns of zero trace. For all \(s \in \mathbf{SU}(2, \mathbf{C})\) , we define an automorphism \(\rho(s)\) of E by \(\rho(s)(m) = sms^*\) . We identify the element
\end{enumerate}

\[
\left( \begin{array}{c c} x & y + i z \\ y - i z & - x \end{array} \right)
\]

of E with the element \((x,y,z)\) of \(R^{3}\) . By imitating the method of Exercise 3, show that \(\rho\) is a morphism of the real Lie group \(\mathbf{SU}(2,\mathbf{C})\) onto the real Lie group \(\mathbf{SO}(3,\mathbf{R})\) with kernel \(\{I,-I\}\) and that \(\mathbf{SO}(3,\mathbf{R})\) is isomorphic to \(\mathbf{SU}(2,\mathbf{C})/\{I,-I\}\)

¶ 6. Let F be the vector space of Hermitian complex matrices with 2 rows and 2 columns. For all \(s \in \mathbf{SL}(2, \mathbf{C})\) , we define an automorphism \(\sigma(s)\) of F by \(\sigma(s)(m) = sms^*\) . We identify the element \(\begin{pmatrix} t + x & y + iz \\ y - iz & t - x \end{pmatrix}\) of F with the element \((t, x, y, z)\) of \(\mathbf{R}^4\) and write

\[
q (t, x, y, z) = t ^ {2} - x ^ {2} - y ^ {2} - z ^ {2}.
\]

Let C be the set of \((t, x, y, z) \in \mathbf{R}^+\) such that \(t^2 - x^2 - y^2 - z^2 > 0\) , \(t > 0\) . Let \(\mathbf{SO}^+(q)\) be the set of \(g \in \mathbf{SO}(q)\) such that \(g(\mathbf{C}) = \mathbf{C}\) ; it is an open normal subgroup of \(\mathbf{SO}(q)\) of index 2 in \(\mathbf{SO}(q)\) (cf.~Exercise 4). Show that \(\sigma\) is a morphism of the underlying real Lie group of \(\mathbf{SL}(2, \mathbf{C})\) onto the real Lie group \(\mathbf{SO}^+(q)\) with kernel \(\{I, -I\}\) . (For the surjectivity, use Exercise 5.) Deduce that the real Lie group \(\mathbf{SO}^+(q)\) is isomorphic to the underlying real Lie group of \(\mathbf{SL}(2, \mathbf{C}) / \{I, -I\}\) , that is (Exercise 3) the underlying real Lie group of \(\mathbf{SO}(3, \mathbf{C})\) .

\begin{enumerate}
\def\labelenumi{\arabic{enumi}.}
\setcounter{enumi}{6}
\tightlist
\item
  \begin{enumerate}
  \def\labelenumii{(\alph{enumii})}
  \tightlist
  \item
    Let G be the set of quaternions of norm 1. This is a Lie subgroup of the real Lie group \(\mathbf{H}^*\) , homeomorphic to \(\mathbf{S}_3\) and hence simply connected, cf.~General Topology, Chapter XI.
  \end{enumerate}
\end{enumerate}

\begin{enumerate}
\def\labelenumi{(\alph{enumi})}
\setcounter{enumi}{1}
\tightlist
\item
  Let E be the real vector space of pure quaternions, identified with \(\mathbf{R}^3\) by
\end{enumerate}

\[
(x, y, z) \mapsto x i + y j + z k.
\]

For all \(g \in \mathbf{G}\) , we define an automorphism \(\rho(g)\) of \(\mathbf{E}\) by \(\rho(g)q = ggg^{-1}\) . Show that \(\rho\) is a morphism of the Lie group \(\mathbf{G}\) onto the real Lie group \(\mathbf{SO}(3, \mathbf{R})\) with kernel \(\{1, -1\}\) . Deduce that the Lie group \(\mathbf{SO}(3, \mathbf{R})\) is isomorphic to the Lie group \(G / \{1, -1\}\) .

\begin{enumerate}
\def\labelenumi{(\alph{enumi})}
\setcounter{enumi}{2}
\item
  Let \(\mathbf{C}\) be identified with the subfield \(\mathbf{R} + \mathbf{R}i\) of \(\mathbf{H}\) . The mapping \((\lambda, q) \mapsto q\lambda\) of \(\mathbf{C} \times \mathbf{H}\) into \(\mathbf{H}\) makes \(\mathbf{H}\) into a vector space over \(\mathbf{C}\) which is identified with \(\mathbf{C}^2\) by the choice of basis \((1, j)\) . For all \(g \in \mathbf{G}\) , let \(\sigma(g)\) be the automorphism of this vector space defined by \(\sigma(g)q = gq\) . Show that \(\sigma\) is an isomorphism of the real Lie group G onto the real Lie group \(\mathbf{SU}(2, \mathbf{C})\) . The latter is therefore simply connected by \((a)\) .
\item
  For \((q_{1}, q_{2}) \in \mathbf{G} \times \mathbf{G}\) , let \(\tau(q_{1}, q_{2})\) be the mapping \(q \mapsto q_{1}q\bar{q}_{2}\) of \(\mathbf{H} = \mathbf{R}^{4}\) into itself. Show that \(\tau\) is a morphism of the real Lie group \(\mathrm{G} \times \mathrm{G}\) onto the real Lie group \(\mathbf{SO}(4, \mathbf{R})\) with kernel \(\mathrm{N} = \{(1, 1), (-1, -1)\}\) and hence that \(\mathbf{SO}(4, \mathbf{R})\) is isomorphic to \((\mathrm{G} \times \mathrm{G}) / \mathrm{N}\) .
\end{enumerate}

\begin{enumerate}
\def\labelenumi{\arabic{enumi}.}
\setcounter{enumi}{7}
\tightlist
\item
  Let \(\mathrm{G}\) be a finite-dimensional connected real Lie group and \(\mathbf{M}\) a finite-dimensional connected manifold of class \(\mathrm{C}^{\infty}\) ; let there be given a law of right operation of class \(\mathrm{C}^{\infty}\) of \(\mathrm{G}\) on \(\mathbf{M}\) . For all \(x \in \mathbf{M}\) , let \(\varphi(x)\) be the corresponding orbital mapping.
\end{enumerate}

\begin{enumerate}
\def\labelenumi{(\alph{enumi})}
\item
  For G to operate transitively on M, it is necessary and sufficient that, for all \(x \in M\) , the mapping \(\mathrm{T}_{e}(\varphi(x))\) of L(G) into \(\mathrm{T}_{x}(\mathrm{M})\) be surjective. (To see that the condition is sufficient, observe that then every G-orbit in M is open.)
\item
  Suppose that G operates transitively on M. We identify M with a homogeneous space H\textbar G, where H is the stabilizer of a point \(x_{0}\) of M. The action of G on M is called imprimitive if there exists a closed submanifold V of M of class \(C^{\omega}\) such that \(0 < dim V < dim M\) and for which every transform V.s (where \(s \in G\) ) is either equal to V or disjoint from V. For this to be so, it is necessary and sufficient that there exist a Lie subgroup L of G such that \(H \subset L \subset G\) and \(\dim(H) < \dim(L) < \dim(G)\) . If this is not so, the action of G on M is called primitive; for this it suffices that there exist no subalgebra of L(G) lying between L(G) and L(H) and distinct from L(G) and L(H).
\item
  Under the hypothesis of (b), let \(\mathcal{L}\) be the image of L(G) under the law of infinitesimal operation associated with the action of G. Let \(\mathcal{E}\) be the algebra of functions of class \(C^\infty\) on M and m the maximal ideal of \(\mathcal{E}\) consisting of the functions in \(\mathcal{E}\) which are zero at \(x_0\) ; we denote by \(\mathcal{L}_p\) for \(p = -1, 0, 1, 2, \ldots\) the set of \(x \in \mathcal{L}\) such that \(\theta(\mathbf{X}) f \in m^{p+1}\) for all \(f \in \mathcal{E}\) ; then \(\mathcal{L}_{-1} = \mathcal{L}\) and \(\mathcal{L}_0\) is the image of L(H) under the law of infinitesimal operation. If (U, \(\phi, \mathbf{R}^n\) ) is a chart of M centered on \(x_0\) and \(\mathbf{X}_1, \ldots, \mathbf{X}_n\) are the vector fields on U defined by the chart, the elements of \(\mathcal{L}_p\) are the \(a_1\mathbf{X}_1 + \cdots + a_n\mathbf{X}_n\) with \(a_1, \ldots, a_n\) in \(m^p|\text{U}\) . Show that \([\mathcal{L}_p, \mathcal{L}_q] \subset \mathcal{L}_{p,q}\) (with the convention \(\mathcal{L}_{-2} = \mathcal{L}\) ). If there exists Y \(\in \mathcal{L}_p\) (for \(p \geqslant 0\) ) such that Y \(\notin \mathcal{L}_{p+1}\) , there exists X \(\in\) L such that {[}Y, X{]} \(\notin \mathcal{L}_p\) . The \(\mathcal{L}_p\) are Lie subalgebras of \(\mathcal{L}\) . For \(p \geqslant 0\) , \(\mathcal{L}_p\) is an ideal of \(\mathcal{L}_0\) . Let \(\rho\) be the canonical linear representation of H on T \(_{x_0}\) (M). The Lie algebra of H/(Ker \(\rho\) ) is isomorphic to \(\mathcal{L}_0/\mathcal{L}_1\) .
\item
  Suppose further that the intersection of the \(\mathcal{L}_p\) is \(\{0\}\) . Then there exists a greatest index \(r\) such that \(\mathcal{L}_r \neq \{0\}\) and all the \(\mathcal{L}_f\) of indices \(\leqslant r\) are distinct. For \(p \geqslant 0\) ,
\end{enumerate}

\[
\dim (\mathcal {L} _ {p} / \mathcal {L} _ {p + 1}) \leqslant \binom {n + p} {p + 1}.
\]

\begin{enumerate}
\def\labelenumi{(\alph{enumi})}
\setcounter{enumi}{4}
\tightlist
\item
  The hypothesis of (d) is satisfied if M is analytic and G operates analytically on M.
\end{enumerate}

\begin{enumerate}
\def\labelenumi{\arabic{enumi}.}
\setcounter{enumi}{8}
\tightlist
\item
  In the notation of Proposition 30, HH' may be an open subset of G dense in G and distinct from G (take G = GL(2, R) with H the upper triangular subgroup and H' the lower strict triangular subgroup).
\end{enumerate}

\exercisesfor{4}

\begin{enumerate}
\def\labelenumi{\arabic{enumi}.}
\tightlist
\item
  Let G be a Lie group, H a normal Lie quasi-subgroup of G and \(\pi\) the canonical mapping of G onto G/H. There exists on G/H one and only one Lie group structure with the following property: for a homomorphism \(\theta\) of G/H into a Lie group \(G'\) to be a Lie group morphism, it is necessary and sufficient that \(\theta \circ \pi\) be a Lie group morphism. Moreover, L(G/H) is canonically isomorphic to L(G)/L(H). (Let Q be a Lie group germ such that L(Q) = L(G)/L(H). Show that, by shrinking Q if necessary, Q can be identified with an open neighbourhood of 0 in G/H and that Proposition 18 of §1 can then be applied to G/H.)
\end{enumerate}

¶ 2. Let G and H be two Lie groups and f a Lie group morphism of G into H.

\begin{enumerate}
\def\labelenumi{(\roman{enumi})}
\item
  The kernel \(\mathbf{N}\) of \(f\) is a normal Lie quasi-subgroup of \(\mathbf{G}\) and \(\mathbf{L}(\mathbf{N}) = \operatorname{Ker} \mathbf{L}(f)\) . (Use exponential mappings of \(\mathbf{G}\) and \(\mathbf{H}\) ).
\item
  Let \(g \colon \mathrm{G} / \mathrm{N} \to f(\mathrm{G})\) be the mapping derived from \(f\) when passing to the quotient. If \(f\) and \(\mathrm{L}(f)\) have closed images, and the topology of \(\mathrm{G}\) admits a countable base, then \(f(\mathrm{G})\) is a Lie quasi-subgroup of \(\mathrm{H}\) with Lie algebra \(\operatorname{Im} \mathrm{L}(f)\) and \(g\) is a Lie group isomorphism, where \(\mathrm{G} / \mathrm{N}\) has the structure defined in Exercise 1. (Reduce it by means of (i) to the case where \(\mathrm{N} = \{e\}\) .)
\end{enumerate}

\begin{enumerate}
\def\labelenumi{\arabic{enumi}.}
\setcounter{enumi}{2}
\tightlist
\item
  Let G be a Lie group, U an open neighbourhood of 0 in L(G) and \(\phi\) an analytic mapping of U into G such that \(\phi(0) = e\) and \(T_{0}\phi = \mathrm{Id}_{\mathrm{L}(G)}\) . The following conditions are equivalent:
\end{enumerate}

\begin{enumerate}
\def\labelenumi{(\alph{enumi})}
\item
  \(\phi\) is an exponential mapping;
\item
  There exists \(m \in \mathbf{Z}\) distinct from 0, 1, -1 such that \(\phi(mx) = \phi(x)^m\) in a neighbourhood of 0.
\end{enumerate}

(If condition (b) holds, show that \(\phi^{*}\) satisfies condition (iii) of Theorem 4.)

¶ 4. (a) Let K = R or C or an ultrametric field with residue field of characteristic \textgreater0. Let G be a Lie group over K, U an open neighbourhood of 0 in L(G) and \(\phi: U \to G\) a mapping which is differentiable at 0 and such
that \(\phi(0) = e\) , \(T_0(\phi) = \mathrm{Id}_{L(G)}\) and \(\phi((\lambda + \lambda')b) = \phi(\lambda b)\phi(\lambda'b)\) for \(\lambda b\) , \(\lambda'b\) , \((\lambda + \lambda')b\) in U. Then \(\phi\) coincides on a neighbourhood of 0 with an exponential mapping.

\begin{enumerate}
\def\labelenumi{(\alph{enumi})}
\setcounter{enumi}{1}
\tightlist
\item
  Let \(k\) be a field of characteristic 0. Let \(\mathbf{K}\) be the valued field \(k((\mathbf{X}))\) , which is of characteristic 0. Let \(\mathbf{G}\) be the additive Lie group \(k[[\mathbf{X}]]\) . Let \(\phi\) be the continuous \(k\) -linear mapping of \(\mathbf{G}\) into \(\mathbf{G}\) such that \(\phi(\mathbf{X}^n) = \mathbf{X}^n + \mathbf{X}^{2n}\) for every integer \(n \geqslant 0\) . Then \(\phi\) satisfies the conditions of (a), but coincides on no neighbourhood of 0 with an exponential mapping.
\end{enumerate}

\begin{enumerate}
\def\labelenumi{\arabic{enumi}.}
\setcounter{enumi}{4}
\tightlist
\item
  Let \((e_1, e_2, e_3)\) be the canonical basis of \(\mathbf{R}^3\) . We give \(\mathbf{R}^3\) the nilpotent Lie algebra structure \(\mathfrak{g}\) such that \([e_1, e_2] = e_3\) , \([e_1, e_3] = [e_2, e_3] = 0\) . Let \(c_{ijk}\) denote the constants of structure of \(\mathfrak{g}\) relative to the basis \((e_1, e_2, e_3)\) .
\end{enumerate}

\begin{enumerate}
\def\labelenumi{(\alph{enumi})}
\tightlist
\item
  Show that on \(\mathbf{R}^3\) the differential forms
\end{enumerate}

\[
\omega_ {1} = d x _ {1}, \quad \omega_ {2} = d x _ {2}, \quad \omega_ {3} = - x _ {1} d x _ {2} + d x _ {3}
\]

satisfy the relations

\[
d \omega_ {k} = - \sum_ {i <   j} c _ {i j k} \omega_ {i} \wedge \omega_ {j} \quad (k = 1, 2, 3)
\]

and are linearly independent at each point of \(\mathbf{R}^3\) .

\begin{enumerate}
\def\labelenumi{(\alph{enumi})}
\setcounter{enumi}{1}
\tightlist
\item
  Deduce that there exists, on an open neighbourhood of 0 in \(\mathbf{R}^3\) , a Lie group germ structure \(\mathbf{G}\) such that \(\mathrm{L}(\mathrm{G}) = \mathfrak{g}\) and
\end{enumerate}

\[
(a _ {1}, a _ {2}, a _ {3}) (x _ {1}, x _ {2}, x _ {3}) = (x _ {1} + a _ {1}, x _ {2} + a _ {2}, x _ {3} + a _ {3} + a _ {1} x _ {2})\tag{1}
\]

for \((a_{1}, a_{2}, a_{3})\) , \((x_{1}, x_{2}, x_{3})\) sufficiently close to 0. Show that in fact formulae (1) define a nilpotent Lie group structure on \(\mathbf{R}^{3}\) .

¶ 6. Let G be a finite-dimensional Lie group, g its Lie algebra, g* the dual vector space of g, σ the adjoint representation of G on g and ρ the representation of G on g* defined by ρ(g) = tσ(g)-1 for all g ∈ G.

\begin{enumerate}
\def\labelenumi{(\alph{enumi})}
\item
  \(\mathrm{L}(\rho)(x) = -^{t}(\mathrm{ad} x)\) for all \(x \in g\) .
\item
  Let \(f \in \mathfrak{g}^*\) . Let \(G_f\) denote the stabilizer of \(f\) in \(G\) ; it is a Lie subgroup of \(G\) and \(\mathfrak{g}_f = L(G_f)\) is the set of \(x \in G\) such that \({}^t (\operatorname{ad}x)f = 0\) .
\item
  For \(x, y\) in \(\mathfrak{g}\) , we write \(\mathrm{B}_f(x, y) = f([x, y])\) . Show that \(\mathrm{B}_f\) is an alternating bilinear form on \(\mathfrak{g}\) , that the mapping \(s_f\) of \(\mathfrak{g}\) into \(\mathfrak{g}^*\) left associated with \(\mathrm{B}_f\) is the mapping \(x \mapsto {}^t (\mathrm{ad}x)\) \(f\) and that the orthogonal of \(\mathfrak{g}\) with respect to \(\mathrm{B}_f\) is \(\mathfrak{g}_f\) . Let \(\beta_f\) denote the non-degenerate alternating bilinear form on \(\mathfrak{g} / \mathfrak{g}_f\) (which is therefore of even dimension) derived from \(\mathrm{B}_f\) when passing to the quotient. We denote by \(\beta_f'\) the inverse form of \(\mathrm{B}_f\) on \((\mathfrak{g} / \mathfrak{g}_f)*\) .
\item
  Let \(\Omega\) denote an orbit of \(G\) on \(\mathfrak{g}^*\) ; we give it the manifold structure derived from that on \(G / G_f\) by transport of structure, where \(f\) is an arbitrary point of \(\Omega\) . Then \(G\) operates analytically on \(\Omega\) on the left. Let \(D\) be the associated law of infinitesimal operation. For all \(a \in \mathfrak{g}\) and all \(f \in \Omega\) ,
\end{enumerate}

\[
\mathrm{D} _ {a} (f) = - (^ {t} \mathrm{ad} a) f.
\]

The tangent subspace \(\mathrm{T}_{f}(\Omega)\) is the image of \(s_{f}\) , that is the orthogonal of \(g_{f}\) in \(g^{*}\) , and is canonically identified with \((\mathfrak{g}/\mathfrak{g}_{f})^{*}\) , so that \(s_{f}\) defines a canonical isomorphism \(r_{f}\) of \(\mathrm{T}_{f}(\Omega)\) onto its dual. The field \(f \mapsto \beta_{f}'\) is an analytic differential form \(\omega\) of degree 2 on \(\Omega\) , which is invariant under G. For \(f \in \Omega\) , \(a \in g\) , \(b \in g\) , \(\omega(\mathrm{D}_{a}(f), \mathrm{D}_{b}(f)) = f([a, b])\) .

\begin{enumerate}
\def\labelenumi{(\alph{enumi})}
\setcounter{enumi}{4}
\item
  Show that \(d\omega = 0\) .
\item
  If \(\alpha\) is a differential form of degree 1 on \(\Omega\) , the isomorphisms \(r_f\) allow us to identify \(\alpha\) with a vector field. Let A be the set of analytic functions on \(\Omega\) with values in K. For \(\phi, \psi\) in A, we write \([\phi, \psi] = \omega(d\phi, d\psi) \in A\) . Show that A has thus a Lie algebra structure.
\item
  For all \(a \in \mathfrak{g}\) , let \(\psi_{a}\) be the function \(f \mapsto f(a)\) on \(\Omega\) . Show that \(a \mapsto \psi_{a}\) is a homomorphism of \(\mathfrak{g}\) into the Lie algebra A and that the form \(d\psi_{a}\) is identified, by means of the isomorphisms \(r_{f}\) , with the vector field \(\mathrm{D}_{a}\) .
\item
  Let U be an open subset of \(\mathfrak{g}^*\) and \(\phi\) an analytic function on U. For all \(f\in \mathrm{U}\) , the differential \(d_{f}\phi\) of \(\phi\) at \(f\) is identified with an element of \(\mathfrak{g}\) . Suppose that \(\mathrm{X}\phi = 0\) for every vector field X defined by the action of G on \(\mathfrak{g}^*\) . Show that then, for all \(f\in \mathrm{U}\) , \(d_{f}\phi\) belongs to the centre of \(\mathfrak{g}_f\) .
\item
  For all \(f \in \mathfrak{g}^*\) , we write \(r_f = \dim \mathfrak{g}_f\) . Let \(r = \inf_{f \in \mathfrak{g}^*} r_f\) . Show that the set \(V\) of \(f \in \mathfrak{g}^*\) such that \(r_f = r\) is open and dense in \(\mathfrak{g}^*\) . Show that, for all \(f \in V\) , \(\mathfrak{g}_f\) is commutative. (Construct \(r\) functions \(\phi\) satisfying the conditions of \((h)\) in a neighbourhood of \(f\) and such that their differentials at \(f\) are linearly independent.)
\end{enumerate}

¶ 7. (a) Let \((\lambda, x) \mapsto \lambda.x\) be a law of continuous left operation of R on T. Then one of the two following holds:

\begin{enumerate}
\def\labelenumi{(\arabic{enumi})}
\item
  either there exists a fixed point and the stabilizer of every point is either \(\mathbf{R}\) or \(\{0\}\) ;
\item
  or there exists a point whose stabilizer is an infinite discrete subgroup of \(\mathbf{R}\) and \(\mathbf{R}\) operates transitively on \(\mathbf{T}\) . (If \(\{\lambda \in \mathbf{R}|\lambda .x = x\} = \{0\}\) , the orbit of \(x\) is homeomorphic to \(\mathbf{R}\) and has frontier points \(\lim_{\lambda \to +\infty}\lambda .x\) which are fixed. If \(\{\lambda \in \mathbf{R}|\lambda .x = x\} = \mathbf{Z}a\) with \(a\neq 0\) , the orbit of \(x\) is homeomorphic to \(\mathbf{T}\) and \(\mathbf{R}\) is transitive.)
\end{enumerate}

\begin{enumerate}
\def\labelenumi{(\alph{enumi})}
\setcounter{enumi}{1}
\item
  If \(f\) is a homeomorphism of \(\mathbf{T}\) onto itself and \(k\) an integer \(\geqslant 1\) , a point \(x \in \mathbf{T}\) is called periodic of period \(k\) for \(f\) if \(f^k(x) = x\) and \(f^h(x) \neq x\) for \(1 \leqslant h < k\) . In the notation of (a), let \(f_{\lambda}\) be the homeomorphism \(x \mapsto \lambda .x\) . Then the set of periodic points of period \(k > 1\) for \(f_{\lambda}\) is empty or equal to \(\mathbf{T}\) . (If there exists a periodic point of period \(k > 1\) for \(f_{\lambda}\) , its stabilizer is distinct from \{0\} and \(\mathbf{R}\) , hence \(\mathbf{R}\) operates transitively on \(\mathbf{T}\) and all the points of \(\mathbf{T}\) are periodic of period \(k\) for \(f_{\lambda}.\) )
\item
  Let \(\Omega\) be a neighbourhood of \(e\) in \(\mathrm{Diff}^{\infty}(\mathbf{T})\) (Differentiable and Analytic Manifolds, R, 15.3.8, Note (1)). There exist an integer \(k > 1\) and \(f\in \Omega\) without fixed point such that the set of periodic points of period \(k\) for \(f\) is neither empty nor equal to \(\mathbf{T}\) . (Let \(p\colon \mathbf{R}\to \mathbf{T}\) be the canonical mapping. Let \(\rho\) be the mapping \(p(x) \mapsto p\left(x + \frac{2\pi}{k}\right)\) of \(\mathbf{T}\) into \(\mathbf{T}\) . Let \(\Omega'\) be a neighbourhood of \(e\) in \(\text{Diff}^\infty(\mathbf{T})\) such that \(\Omega'^2 \subset \Omega\) . For sufficiently large \(k, \rho \in \Omega'\) . On the other hand, let \(\sigma\) be a mapping of \(\mathbf{T}\) into \(\mathbf{T}\) such that \(\sigma(y) = y\) for
\end{enumerate}

\[
y \notin p \left(\right] 0, \frac {2 \pi}{k} [)
\]

and \(\sigma(p(x)) = p(x + \lambda(x))\) for \(x \in \left]0, \frac{2\pi}{k}\right\}\) , where

\[
0 <   \lambda (x) <   \frac {2 \pi}{k} - x;
\]

can be chosen such that \(\sigma \in \Omega'\) . Then \(f = \rho \circ \sigma\) has the required properties.) Such an \(f\) cannot, by (b), be in the image of a continuous morphism of \(\mathbf{R}\) into \(\operatorname{Diff}^{\infty}(\mathbf{T})\) .

\begin{enumerate}
\def\labelenumi{(\alph{enumi})}
\setcounter{enumi}{3}
\item
  Let Y be a compact differentiable manifold of dimension \(\geqslant1\) . Every neighbourhood of e in \(\operatorname{Diff}^{\infty}(Y)\) contains an element h with the following property: h belongs to the image of no continuous morphism of R into \(\operatorname{Diff}^{0}(Y)\) . (The manifold Y contains an open submanifold of the form \(T \times D\) , where D is the open Eulcidean ball of radius 1 and centre 0 in \(R^{n}\) ( \(n = \dim Y - 1\) ). Let \(g \in \operatorname{Diff}^{\infty}(D)\) be such that \(g(0) = 0\) , \(\|g(x)\| > \|x\|\) for \(0 < \|x\| < \frac{1}{2}\) and \(g(x) = x\) for \(\|x\| \geqslant \frac{1}{2}\) and g very close to e in \(\operatorname{Diff}^{\infty}(D)\) . Let f be as in (e), and, for 0 \textless{} t \textless{} 1, let \(f_{i}: T \to T\) be defined as follows: for \(x \in T\) , \(y \in p^{-1}(x)\) , \(z \in p^{-1}(f(x))\) and \(|z - y| < \pi\) , we write \(f_{i}(x) = p(ty + (1 - t)z)\) . Then there exists \(h \in \operatorname{Diff}^{\infty}(Y)\) such that \(h(x, y) = (f_{2\|x\|}(x), g(y))\) for \(x \in T\) , \(y \in D\) and \(h(u) = u\) for \(u \in Y - (T \times D)\) ; moreover, we can choose f and g such that h is arbitrarily close to e in \(\operatorname{Diff}^{\infty}(Y)\) . The points of \(T \times \{0\}\) are the \(y \in Y\) such that when n tends to +∞, \(h^{kn}(x)\) tends to a point not fixed under h (in fact it is periodic of period k). Therefore, every homeomorphism of Y onto Y which commutes with h leaves \(T \times \{0\}\) invariant. Then apply (e).
\item
  Let Y be a compact differentiable manifold of dimension \(\geqslant1\) . There do not exist a Lie group G and continuous morphisms \(\operatorname{Diff}^{\infty}(Y)\to G\) , \(G\to\operatorname{Diff}^{0}(Y)\) whose composition is the canonical injection of \(\operatorname{Diff}^{\infty}(Y)\) into \(\operatorname{Diff}^{0}(Y)\) . (Use (d).)
\end{enumerate}

\begin{enumerate}
\def\labelenumi{\arabic{enumi}.}
\setcounter{enumi}{7}
\tightlist
\item
  Let G be a finite-dimensional Lie group. Suppose that L(G) is simple. Let A be a normal subgroup of G. If A is not open, A is discrete and its centralizer in G is open. (Let a be the subalgebra tangent to A at e. Show that \(a = \{0\}\) . Let x be an element of A not belonging to the centre of G. Consider the mapping \(y \mapsto xyj^{-1}\) of G into A. Deduce from the relation \(a = \{0\}\) that the centralizer of x in G is open. Then, using an exponential mapping, show that the centralizer of x in G contains a neighbourhood of e independent of x.)
\end{enumerate}

¶ 9. Let G be a compact Lie group over K of dimension n.~Suppose that the Lie algebra L(G) is simple. For all \(g \in G\) , every integer \(m \geqslant 1\) and every neighbourhood V of e in G, let M(g, m, V) denote the set of elements of G of the form \(\prod_{i=1} x_i(y_i, g) x_i^{-1}\) with \(x_1, \ldots, x_m, y_1, \ldots, y_m\) in V.

\begin{enumerate}
\def\labelenumi{(\alph{enumi})}
\tightlist
\item
  Show that if \(g\) is an element of \(G\) whose centralizer is not open, if \(m\) is an integer \(\geqslant n\) and if \(V\) is a neighbourhood of \(e\) , then \(M(g, m, V)\) is a neighbourhood of \(e\) . (There exists \(a \in L(G)\) such that \((\operatorname{Ad} g)(a) \neq a\) . Let \(\phi\) be an exponential mapping of \(G\) . Let \(y\) be the image of 1 under the tangent mapping at 0 to \(\lambda \mapsto \langle \phi(xa), g \rangle\) (where \(\lambda \in K\) ). Then \(b = (\operatorname{Ad} g - 1)(a) \neq 0\) . There exist \(x_1, \ldots, x_m\) in \(V\) such that \(\{(\operatorname{Ad} x_1)(b), \ldots, (\operatorname{Ad} x_m)(b)\}\) contains a basis of \(L(G)\) . Consider the mapping
\end{enumerate}

\[
\left(x _ {1} ^ {\prime}, \dots , x _ {m} ^ {\prime}, y _ {1}, \dots , y _ {m}\right) \mapsto \prod_ {i = 1} ^ {n} x _ {i} ^ {\prime} (y _ {i}, g) x _ {i} ^ {\prime - 1}
\]

of \(\mathbf{G}^{2m}\) into G. Show that the tangent mapping at \((x_{1},\ldots ,x_{m},e,\ldots ,e)\) is surjective.)

\begin{enumerate}
\def\labelenumi{(\alph{enumi})}
\setcounter{enumi}{1}
\item
  Let \(\mathbf{U}\) be a neighbourhood of \(e\) in \(\mathbf{G}\) . Let \(m\) be an integer \(\geqslant 1\) . There exists an element \(g\) of \(\mathbf{G}\) whose centralizer is not open, such that \(\mathbf{M}(g, m, \mathbf{G}) \subset \mathbf{U}\) . (Argue by reductio ad absurdum using the compactness of \(\mathbf{G}\) .)
\item
  Let \(\mathbf{G}'\) be a compact Lie group over \(\mathbf{K}\) whose Lie algebra is simple. Let \(\phi: \mathbf{G} \to \mathbf{G}'\) be a bijective homomorphism of abstract groups. Then \(\phi\) is continuous. (Apply (b) to \(\mathbf{G}'\) and (a) to \(\mathbf{G}\) .)
\end{enumerate}

\begin{enumerate}
\def\labelenumi{\arabic{enumi}.}
\setcounter{enumi}{9}
\item
  Let \(G\) be a Lie group over \(K\) . Show that there exists a neighbourhood \(V\) of \(e\) with the following property: for every sequence \((x_{n})\) of elements of \(V\) , if we define inductively \(y_{1} = x_{1}, y_{n} = (x_{n}, y_{n-1})\) , the sequence \((y_{n})\) tends to \(e\) .
\item
  \begin{enumerate}
  \def\labelenumii{(\alph{enumii})}
  \tightlist
  \item
    For \(x, y\) in \(\mathbf{S}_{2n-1} \subset \mathbf{C}^n\) , we define \(\alpha(x, y) \in [0, \pi]\) by
  \end{enumerate}
\end{enumerate}

\[
\cos \alpha (x, y) = \mathscr {R} ((x | y)).
\]

For \(s, t\) in \(\mathbf{U}(n)\) , we write \(d(s, t) = \sup_{x \in \mathbf{S}_{2n-1}} \alpha(sx, tx)\) . Show that \(d\) is a left and right invariant distance on \(\mathbf{U}(n)\) .

\begin{enumerate}
\def\labelenumi{(\alph{enumi})}
\setcounter{enumi}{1}
\item
  Let \(s \in \mathbf{U}(n)\) . Let \(\theta_1, \ldots, \theta_m\) be the numbers in \(\mathbf{J} - \pi, \pi \mathbf{j}\) such that the \(e^{i\theta_j}\) are the distinct eigenvalues of \(s\) . Let \(\theta(s) = \sup_{1 \leqslant j \leqslant m} |\theta_j|\) . Then \(d(e, s) = \theta(s)\) . (Let \(V_j\) be the eigenspace of \(s\) corresponding to \(e^{i\theta_j}\) . For \(x \in S_{2n-1}\) , find a lower bound to \(\mathcal{R}((x|sx))\) by decomposing with respect to the \(V_j\) .)
\item
  Let \(s, t\) be in \(\mathbf{U}(n)\) be such that \(\theta(t) < \frac{\pi}{2}\) and \(s\) commutes with \((s, t)\) . Then \(s\) and \(t\) commute. (In the notation of (b), let \(\mathrm{V}_j'\) be the orthogonal supplement of \(\mathrm{V}_j\) and \(\mathrm{W}_j = t(\mathrm{V}_j)\) ; show that \(\mathrm{W}_j = (\mathrm{W}_j \cap \mathrm{V}_j) + (\mathrm{W}_j \cap \mathrm{V}_j')\) and then that \(\mathrm{W}_j \cap \mathrm{V}_j' = \{0\}\) .)
\item
  Show that there exists a compact neighbourhood \(\mathbf{V}\) of \(e\) in \(\mathbf{U}(n)\) which has the property of Exercise 10, which is stable under inner automorphisms of \(\mathbf{U}(n)\) and which is such that if \(x, y\) are non-permutable elements of \(\mathbf{V}\) then \(x\) and \((x, y)\) are not permutable. (Use (c).)
\item
  Let \(\beta\) be the normalized Haar measure of \(\mathbf{U}(n)\) . Let \(\mathrm{V}_1\) be a symmetric compact neighbourhood of \(e\) in \(\mathbf{U}(n)\) such that \(\mathrm{V}_1^2 \in \mathrm{V}\) . Let \(p\) be an integer such that \(\beta (\mathrm{V}_1) > 1 / p\) . Show that, for every finite subgroup F of \(\mathbf{GL}(n,\mathbf{C})\) , there exists a commutative normal subgroup A of F such that \(\operatorname{Card}(\mathrm{F} / \mathrm{A}) \leqslant p\) (Jordan's theorem). (Reduce it to the case where \(\mathrm{F} \subset \mathbf{U}(n)\) . Take A to be the subgroup of F generated by \(\mathrm{F} \cap \mathrm{V}\) and use (d).)
\end{enumerate}

\begin{enumerate}
\def\labelenumi{\arabic{enumi}.}
\setcounter{enumi}{11}
\tightlist
\item
  Let G be a finite-dimensional connected real or complex Lie group and \(\alpha\) a left invariant differential form of degree p on G. For \(\alpha\) to be right invariant too, it is necessary and sufficient that \(d\alpha = 0\) . (Observe that the condition for right invariance can be written
\end{enumerate}

\[
\wedge^ {p} (^ {t} \mathrm{Ad} (s)). \alpha (e) = \alpha (e)
\]

for all \(s \in G\) and hence it is equivalent to the condition

\[
\sum_ {j = 1} ^ {f} \left\langle \alpha (e), u _ {1} \wedge \dots \wedge u _ {j - 1} \wedge [ u, u _ {j} ] \wedge u _ {j + 1} \wedge \dots \wedge u _ {p} \right\rangle = 0
\]

for all \(u, u_1, \ldots, u_p\) in \(\mathrm{L}(\mathrm{G})\) . In particular, the left and right invariant differential forms of degree 1 on \(\mathbf{G}\) form a vector space of dimension

\[
\dim \mathrm{L} (\mathrm{G}) - \dim [ \mathrm{L} (\mathrm{G}), \mathrm{L} (\mathrm{G}) ].
\]

\begin{enumerate}
\def\labelenumi{\arabic{enumi}.}
\setcounter{enumi}{12}
\tightlist
\item
  Let \(\mathbf{R}^n\) be given the usual scalar product \(((\xi_i), (\eta_i)) \mapsto \sum_{i=1} \xi_i \eta_i\) . Let \(\mathbf{I}(n)\) be the group of isometric affine transformations of \(\mathbf{R}^n\) . It is canonically isomorphic to the Lie subgroup of \(\mathbf{GL}(n + 1, \mathbf{R})\) consisting of the matrices
\end{enumerate}

\[
S = \left( \begin{array}{c c} U & x \\ 0 & 1 \end{array} \right)
\]

where \(U \in \mathbf{O}(n)\) and \(x\) is an arbitrary element of \(\mathbf{R}^n\) (matrix of type \((n,1)\) ). \(\mathbf{I}(n)\) can be identified with the semi-direct product of \(\mathbf{O}(n)\) by \(\mathbf{R}^n\) defined by the canonical injection of \(\mathbf{O}(n)\) into \(\mathbf{GL}(n,\mathbf{R})\) (the matrix \(S\) is then denoted by \((U,x)\) ). Let \(p\colon \mathbf{I}(n)\to \mathbf{O}(n)\) be the canonical surjection, with kernel \(\mathbf{R}^n\) .

\begin{enumerate}
\def\labelenumi{(\alph{enumi})}
\tightlist
\item
  Let \(\Gamma\) be a discrete subgroup of \(\mathrm{I}(n)\) . Consider the closure \(\tilde{p}(\Gamma)\) of \(p(\Gamma)\) in \(\mathbf{O}(n)\) . Show that its identity component is commutative. (Let \(V\) be a neighbourhood of \(I\) in \(\mathbf{O}(n)\) with the properties of Exercise 11 ( \(d\) ) and such that further \(\|U - I\| \leqslant \frac{1}{4}\) for all \(U \in V\) (the norm on \(\operatorname{End}(\mathbf{R}^n)\) being derived from the Euclidean norm on \(\mathbf{R}^n\) ). Argue by reductio ad absurdum, assuming that there exist \(S_1 = (U_1, x_1)\) , \(S_2 = (U_2, x_2)\) such that \(U_1\) and \(U_2\) belong to \(V\) and do not commute. Show that \((S_1, S_2) = ((U_1, U_2), y)\) with
\end{enumerate}

\[
\| y \| \leqslant \frac {1}{4} (\| x _ {1} \| + \| x _ {2} \|).
\]

Then define the \(S_{k}\) inductively by \(S_{k}=(S_{1},S_{k-1})\) for \(k\geqslant3\) . Using the choice of V, show that this sequence has an infinity of distinct terms and is bounded in \(\mathbf{I}(n)\) , which is absurd.)

\begin{enumerate}
\def\labelenumi{(\alph{enumi})}
\setcounter{enumi}{1}
\item
  \(\Gamma\) is called crystallographic if \(\mathrm{I}(n)/\Gamma\) is compact. Show that if this is so, for all \(x \in \mathbf{R}^n\) , the affine subspace \(\mathbf{L}\) of \(\mathbf{R}^n\) generated by \(\Gamma x\) is equal to \(\mathbf{R}^n\) . (Otherwise, for all \(y \in \mathbf{R}^n\) , all the points of \(\Gamma y\) are the same distance from \(\mathbf{L}\) and, as this distance can be arbitrarily large \(\mathrm{I}(n)/\Gamma\) is not compact.)
\item
  If \(\Gamma\) is commutative and crystallographic, then \(\Gamma \subset \mathbf{R}^n\) . (If \(S = (U, x) \in \Gamma\) is such that the vector subspace \(\mathrm{V} \subset \mathbf{R}^n\) consisting of the points invariant under \(U\) is not equal to \(\mathbf{R}^n\) and \(\mathrm{V}^\perp\) denotes the orthogonal subspace of \(\mathrm{V}\) , then, for all \(S' = (U', x') \in \Gamma\) , \(U'\) leaves \(\mathrm{V}\) and \(\mathrm{V}^\perp\) stable. On the other hand, since \(U|\mathrm{V}^\perp\) has no non-zero invariant point, show that this allows us to assume, by changing the origin, that \(x \in \mathrm{V}\) . Using (b), there exists \(S' = (U', x') \in \Gamma\) such that the orthogonal projection \(y'\) of \(x'\) onto \(\mathrm{V}^\perp\) is \(\neq 0\) . By evaluating \(Sx'\) by the formula \(S = S'SS'^{-1}\) , deduce that \(U.y' = y'\) contrary to the definition of \(\mathrm{V}\) .
\item
  If \(\Gamma\) is crystallographic, \(\Gamma \cap \mathbf{R}^n\) is a free commutative group of rank \(n\) and \(\Gamma / (\Gamma \cap \mathbf{R}^m)\) is a finite group (Bieberbach's Theorem). (Let \(W\) be the vector subspace of \(\mathbf{R}^n\) generated by \(\Gamma \cap \mathbf{R}^n\) . The compact group \(p(\Gamma)\) has only a finite number of connected components. If \(W = \{0\}\) , \(\Gamma\) contains, by \((a)\) , a commutative normal subgroup of finite index \(\Gamma_1\) ; \(\Gamma_1\) is then crystallographic, which implies a contradiction with \((c)\) . Hence \(W \neq \{0\}\) . Show that \(p(\Gamma)\) leaves \(W\) stable and that \(p(\Gamma)|W\) is finite; otherwise there would exist a sequence \((S_m)\) in \(\Gamma\) such that if \(S_m = (U_m, x_m)\) , the \(U_m\) are distinct and tend to \(I\) ; then form the \((I, a_j)S_m(I, a_j)^{-1s_m^1}\) , where \((a_j)\) is a basis of \(\Gamma \cap \mathbf{R}^n\) , and obtain a contradiction with the hypothesis that \(\Gamma\) is discrete. Finally, to see that \(W = \mathbf{R}^n\) , show that otherwise the action of \(\Gamma\) on \(\mathbf{R}^n/W\) would be that of a crystallographic group containing no translation \(\neq 0\) .)
\end{enumerate}

\exercisesfor{5}

\begin{enumerate}
\def\labelenumi{\arabic{enumi}.}
\tightlist
\item
  Suppose that \(\mathbf{K} = \mathbf{R}\) . Let \(r\) be an integer \(\geqslant 1\) or \(\infty\) . A group of class \(\mathbf{C}^r\) is a set \(\mathbf{G}\) with a group structure and a manifold structure of class \(\mathbf{C}^r\) such that the mapping \((x,y)\mapsto xy^{-1}\) of \(\mathbf{G}\times \mathbf{G}\) into \(\mathbf{G}\) is of class \(\mathbf{C}^r\) .
\end{enumerate}

\begin{enumerate}
\def\labelenumi{(\alph{enumi})}
\tightlist
\item
  Suppose henceforth that \(r \geqslant 2\) . A neighbourhood of \(e\) in \(G\) is identified with a neighbourhood of 0 in a Banach space. Let \(xy = P(x, y)\) , where \(P\) is a mapping of class \(C^2\) on an open neighbourhood of \((0, 0)\) . Let
\end{enumerate}

\[
\left(\mathrm{D} _ {1} \mathrm{D} _ {2} \mathrm{P}\right) (0, 0) = \mathrm{B}.
\]

Show that

\[
x y = x + y + \mathrm{B} (x, y) + | x | | y | o (1)
\]

as \((x,y)\to (0,0)\) . (Use an expansion of P of order 2 with integral remainder.)

\begin{enumerate}
\def\labelenumi{(\alph{enumi})}
\setcounter{enumi}{1}
\tightlist
\item
  Show that
\end{enumerate}

\[
\begin{array}{c} x ^ {- 1} = - x + \mathrm{B} (x, x) + | x | ^ {2} o (1) \\ x y x ^ {- 1} = y + \mathrm{B} (x, y) - \mathrm{B} (y, x) + | x | | y | o (1) \\ x ^ {- 1} y ^ {- 1} x y = \mathrm{B} (x, y) - \mathrm{B} (y, x) + | x | | y | o (1) \end{array}
\]

as \((x,y)\) tends to \((0,0)\) .

\begin{enumerate}
\def\labelenumi{\arabic{enumi}.}
\setcounter{enumi}{1}
\tightlist
\item
  Let \(G\) be a group of class \(C^r\) with \(r \geqslant 2\) .
\end{enumerate}

\begin{enumerate}
\def\labelenumi{(\alph{enumi})}
\item
  Let \(t \in \mathcal{T}^{(s)}(\mathrm{G})\) , \(t' \in \mathcal{T}^{(s')}(\mathrm{G})\) . If \(s + s' \leqslant r\) , define \(t * t' \in \mathcal{T}^{(s + s')}(\mathrm{G})\) as in § 3, no. 1. If \(t\) and \(t'\) are without constant term, \(t * t'\) is without constant term. The image of \(t\) under the mapping \(x \mapsto x^{-1}\) of \(\mathrm{G}\) into \(\mathrm{G}\) is denoted by \(t^{\vee}\) ; then \(t^{\vee} \in \mathcal{T}^{(s)}(\mathrm{G})\) . If \(f\) is a function of class \(\mathrm{C}^r\) on \(\mathrm{G}\) with values in a Hausdorff polynormed space, \(f * t\) denotes the function on \(\mathrm{G}\) defined by \((f * t)(x) = \langle \varepsilon_x * t'^{\nu}, f \rangle\) for all \(x \in \mathrm{G}\) ; this function is of class \(\mathrm{C}^r - s\) if \(s < \infty\) . Show, as in § 3, no. 4, that \(f * (t * t') = (f * t) * t'\) .
\item
  If \(t \in \mathrm{T}_e(\mathrm{G})\) , the vector field \(x \mapsto \varepsilon_x * t\) on \(\mathbf{G}\) is denoted by \(\mathbf{L}_t\) . Show, as in § 3, no. 6, that, for \(t, t'\) in \(\mathrm{T}_e(\mathrm{G})\) , \(\mathrm{L}_{t^* t'} = \mathrm{L}_t \circ \mathrm{L}_{t'}\) and hence that
\end{enumerate}

\[
\mathrm{L} _ {t ^ {*} t ^ {\prime} - t ^ {*} t} = \left[ \mathrm{L} _ {t}, \mathrm{L} _ {t ^ {\prime}} \right].
\]

\[
t * t ^ {\prime} - t ^ {\prime} * t
\]

\[
\mathrm{T} _ {e} (\mathrm{G})
\]

\begin{enumerate}
\def\labelenumi{(\alph{enumi})}
\setcounter{enumi}{2}
\tightlist
\item
  Using the notation of Exercise 1 (a), if \(t \in \mathrm{T}_e(\mathrm{G})\) , then
\end{enumerate}

\[
[ t, t ^ {\prime} ]
\]

\[
\mathrm{L} _ {t} (x) = (\mathrm{D} _ {2} \mathrm{P} (x, 0)) (t)
\]

for \(x\) sufficiently close to 0. Deduce that \([t, t'] = \mathbf{B}(t, t') - \mathbf{B}(t', t)\) .

\begin{enumerate}
\def\labelenumi{(\alph{enumi})}
\setcounter{enumi}{3}
\tightlist
\item
  Show that \(\mathrm{T}_e(\mathrm{G})\) , with the bracket \((t, t') \mapsto [t, t']\) , is a normable Lie algebra. (To prove the Jacobi identity, use (c), Exercise 1 (b) and identity (5) of Algebra, Chapter I, §6, no. 2.)
\end{enumerate}

(For a sequel to this exercise, cf.~§ 8, Exercise 6.)

\exercisesfor{6}

\begin{enumerate}
\def\labelenumi{\arabic{enumi}.}
\item
  Let D be a set of elements of \(\mathbf{SL}(2,\mathbf{C})\) of the form \(\begin{pmatrix}a & b \\ 0 & a^{-1}\end{pmatrix}\) where a \textgreater{} 0 and \(b \in C\) . Show that D is a Lie subgroup of the underlying real Lie group of \(\mathbf{SL}(2,\mathbf{C})\) and that the mapping \((u,d) \mapsto ud\) of \(\mathbf{SU}(2,\mathbf{C}) \times \mathbf{D}\) into \(\mathbf{SL}(2,\mathbf{C})\) is an isomorphism of real analytic manifolds. Deduce from this and §3, Exercise 7 (c) that \(\mathbf{SL}(2,\mathbf{C})\) is simply connected.
\item
  Let \(G\) be the universal covering of \(\mathbf{SL}(2, \mathbf{R})\) , \(\pi\) the canonical morphism of \(G\) onto \(\mathbf{SL}(2, \mathbf{R})\) and \(N\) its kernel.
\end{enumerate}

\begin{enumerate}
\def\labelenumi{(\alph{enumi})}
\item
  Arguing as in Exercise 1, show that there exists an isomorphism of the real analytic manifold \(\mathbf{U} \times \mathbf{R}^2\) onto the real analytic manifold \(\mathbf{SL}(2, \mathbf{R})\) . Deduce that \(N\) is isomorphic to \(\mathbf{Z}\) .
\item
  If \(\rho\) is an analytic linear representation of \(G\) on a complex vector space, then \(\operatorname{Ker} \rho \supset N\) . (The representation \(\mathrm{L}(\rho)\) of \(\mathfrak{sl}(2, \mathbf{R})\) defines, by complexification, a representation of \(\mathfrak{sl}(2, \mathbf{C})\) and the latter is, by Exercise 1, of the form \(\mathrm{L}(\sigma)\) , where \(\sigma\) is an analytic linear representation of \(\mathbf{SL}(2, \mathbf{C})\) . Then \(\mathrm{L}(\sigma \circ \pi) = \mathrm{L}(\rho)\) and hence \(\sigma \circ \pi = \rho\) .)
\end{enumerate}

\begin{enumerate}
\def\labelenumi{\arabic{enumi}.}
\setcounter{enumi}{2}
\item
  Let \(G_{1}\) be the simply connected nilpotent real Lie group defined in Exercise 5 of § 4. Let Z be the centre of \(G_{1}\) , N a non-trivial discrete subgroup of Z and \(G = G_{1}/N\) . If \(\rho\) is a finite-dimensional analytic linear representation of G, \(\rho(Z/N)\) is a semi-simple set of automorphisms, for Z/N is compact; but \(Z = (G, G)\) and hence \(\rho(Z/N)\) is unipotent (Chapter I, § 6, Proposition 6). Therefore \(\rho\) is trivial on Z/N.
\item
  Let G be a compact connected complex Lie group. Show that every analytic linear representation of G is trivial.
\item
  In the notation of Exercise 9 of § 1, show that \(\mathbf{H}'\) is an integral subgroup of H. Deduce that in the simply connected group \(\mathbf{SU}(2,\mathbf{C})\times \mathbf{SU}(2,\mathbf{C})\) (§ 3, Exercise 7 (c)), there exist one-parameter subgroups which are not closed.
\end{enumerate}

¶ 6. Let I be the interval \(\{1,2\}\) with the discrete topology. Let E be the complete normed space of functions \(f: I \to R\) such that \(\|f\| = \sum_{i \in I} |f(i)| < +\infty\) . For all \(x \in I\) , let \(\varepsilon_x\) be the element of E such that \(\varepsilon_x(t) = 0\) for \(t \neq x\) and \(\varepsilon_x(x) = 1\) . Let P be the subgroup of E generated by the \(x\varepsilon_x\) with \(x \in I\) ; it is a discrete subgroup of E. Let G be the Lie group E/P. Let F be the hyperplane of E consisting of the \(f \in E\) such that \(\sum_{i \in I} f(i) = 0\) . Let H be the Lie group F/(F ∩ P). Let \(\phi\) be the canonical morphism of H into G. Show that \(\phi\) is bijective and is an immersion but is not a Lie group isomorphism. (Let f be a linear form on E with kernel F; show that \(f(P) = R\) and hence that \(F + P = E\) .) Deduce that in Proposition 3 the countability hypothesis cannot be omitted.

\begin{enumerate}
\def\labelenumi{\arabic{enumi}.}
\setcounter{enumi}{6}
\tightlist
\item
  Let \(\phi\) be the morphism of \(\mathbf{Z} / 2\mathbf{Z}\) into the permutation group of \(\mathbf{C}\) which maps the non-identity element to \(z\mapsto \overline{z}\) . Let \(\mathrm{G}\) be the semi-direct product of \(\mathbf{Z} / 2\mathbf{Z}\) by the real Lie group \(\mathbf{C}\) corresponding to \(\phi\) . It is a real Lie group with lie algebra \(\mathbf{R}^2\) . Show that there exists on \(\mathrm{G}\) no complex Lie group structure compatible with the real Lie group structure.
\end{enumerate}

¶ 8. Let H be a complex Hilbert space of dimension \(\aleph_{0}\) and G the unitary group of H considered as a real Lie group. The group G is simply connected. \(\dagger\) (a) Let Z be the centre of G, which is isomorphic to T. Let a be an irrational number. Let \(\hat{\mathbf{a}}\) be the Lie subalgebra of \(\mathrm{L}(\mathrm{G})\times \mathrm{L}(\mathrm{G})\) consisting of the \((x,ax)\) where \(x\in \mathrm{L}(Z)\) . Let S be the corresponding integral subgroup of \(\mathrm{G}\times \mathrm{G}\) . Then \(\hat{\mathbf{a}}\) is an ideal of \(\mathrm{L}(\mathrm{G})\times \mathrm{L}(\mathrm{G})\) ; however S is not closed and is dense in \(\mathbf{Z}\times \mathbf{Z}\) .

\begin{enumerate}
\def\labelenumi{(\alph{enumi})}
\setcounter{enumi}{1}
\tightlist
\item
  Let \(\mathfrak{g} = (\mathrm{L}(\mathrm{G})\times \mathrm{L}(\mathrm{G})) / \tilde{\mathfrak{s}}\) . There exists no Lie group with Lie algebra \(\mathfrak{g}\) . (Let \(\mathrm{H}\) be such a group. As \(\mathrm{G}\) is simply connected, there exists a morphism \(\phi : \mathrm{G}\times \mathrm{G}\to \mathrm{H}\) such that \(\mathrm{L}(\phi)\) is the canonical mapping of \(\mathrm{L}(\mathrm{G})\times \mathrm{L}(\mathrm{G})\) onto \(\mathfrak{g}\) . Let \(\mathrm{N} = \operatorname{Ker}\phi\) . Then \(\mathrm{L}(\mathrm{N})\supset \tilde{\mathfrak{s}}\) , hence \(\mathrm{N}\supset \mathrm{S}\) and therefore \(\mathrm{N}\supset \mathrm{Z}\times \mathrm{Z}\) by (a). Then \(\mathrm{L}(\mathrm{N})\supset \mathrm{L}(\mathrm{Z})\times \mathrm{L}(\mathrm{Z})\) , which is absurd.)
\end{enumerate}

\begin{enumerate}
\def\labelenumi{\arabic{enumi}.}
\setcounter{enumi}{8}
\item
  Let X be a non-empty connected compact complex manifold of dimension n.~Suppose that there exist holomorphic vector fields \(\xi_{1},\ldots,\xi_{n}\) on X which are linearly independent at each point of X. Show that there exist a complex Lie group G and a discrete subgroup D of G such that X is diffeomorphic to G/D. (Let \(c_{ijk}\) be the holomorphic functions on X such that \([\xi_{i},\xi_{j}]=\sum_{k}c_{ijk}\xi_{k}\) . The \(c_{ijk}\) are constants because X is compact. Take G to be the simply connected complex Lie group whose Lie algebra admits the \(c_{ijk}\) as constants of structure and use Theorem 5.)
\item
  Let G be a Lie group with a finite number of connected components. For G to be unimodular, it is necessary and sufficient that Tr ad \(a = 0\) for all \(a \in \mathrm{L}(\mathrm{G})\) .
\item
  Let E be a complete normable space over \(\mathbf{C}\) , \(v \in \mathcal{L}(\mathrm{E})\) and \(g = \exp(v)\) . Suppose that \(\operatorname{Sp}(v) \cap 2i\pi(\mathbf{Z} - \{0\}) = \emptyset\) . Let \(\mathrm{E}_1, \mathrm{E}_2\) be closed vector subspaces of E which are stable under \(v\) , such that \(\mathrm{E}_2 \subset \mathrm{E}_1\) . Suppose that the automorphism of \(\mathrm{E}_1 / \mathrm{E}_2\) derived from \(g\) is the identity. Then \(v(\mathrm{E}_1) \subset \mathrm{E}_2\) .
\end{enumerate}

¶ 12. Consider a real or complex complete normable space E and a closed vector subspace F of E such that \(F^{0}\) admits no topological supplement in \(E^{\prime}\dagger\) (a) Let A (resp. B) be the complete normable space of continuous endomorphisms of E (resp. F). Let C be the set of \(u \in A\) such that \(u(F) \subset F\) . Let \(\alpha\) be the mapping \(u \mapsto (u, u \mid F)\) of C into \(A \times B\) . Then \(\alpha\) is an isomorphism of the complete normable space C onto a closed vector subspace of \(A \times B\) and \(\alpha(C)\) admits no topological supplement in \(A \times B\) . (Let \(x \in E\) be such that \(x \notin F\) . Let \(\xi \in \mathrm{F}^0\) be such that \(\langle x, \xi \rangle = 1\) . If \(\eta \in \mathrm{E}'\) , let \(\tilde{\eta}\) denote the element \(y \mapsto \langle y, \eta \rangle x\) of A. Suppose that there exists a projector \(\pi\) of A × B onto \(\alpha(\mathrm{C})\) . We define \(\tilde{\pi} \colon \mathrm{E}' \to \mathrm{E}'\) by \(\tilde{\pi}(\eta) = t(\alpha^{-1}(\pi(\tilde{\eta}, 0)))(\xi)\) . Then \(\tilde{\pi}\) is a projector of \(\mathrm{E}'\) onto \(\mathrm{F}^0\) , which is absurd.)

as follows: \(M = \bigoplus_{i=0}^{n} M_{i}\) is graded by the \(M_{i}\) ; \(M_{0}\) is the set of scalar multiples of 1; \(M_{i}\) is the set of scalar multiples of a non-zero element u; \(M_{2} = \{0\}\) ; \(M_{3} = F\) ; \(M_{4} = E\) ; \(M_{i} = \{0\}\) for i \textgreater{} 5; ux = x for all \(x \in M_{3}\) . Let N be the complete normable space of continuous endomorphisms of the complete normable space M. Let \(N_{1}\) be the set of continuous derivations of M. Show, using (a), that \(N_{1}\) has no topological supplement in N.

\begin{enumerate}
\def\labelenumi{(\alph{enumi})}
\setcounter{enumi}{2}
\tightlist
\item
  Show that the group of bicontinuous automorphisms of M is a Lie quasi-subgroup of \(\mathbf{GL}(\mathbf{M})\) , but not a Lie subgroup of \(\mathbf{GL}(\mathbf{M})\) . (Use Corollary 2 to Proposition 18 and (b).)
\end{enumerate}

\begin{enumerate}
\def\labelenumi{\arabic{enumi}.}
\setcounter{enumi}{12}
\item
  Let \(\mathbf{G}\) be a real Lie group, \(\mathbf{L} = \mathbf{L}(\mathbf{G})\) , and \(\phi: \mathbf{L} \to \mathbf{G}\) a mapping which is differentiable at 0, such that \(\mathrm{T}_0(\phi) = \mathrm{Id}_{\mathrm{L}}\) and such that \(\phi(nx) = \phi(x)^n\) for all \(x \in \mathbf{L}\) and \(n \in \mathbf{Z}\) . Then \(\phi = \exp_0\) . (Let \(\mathbf{V}\) be a neighbourhood of 0 in \(\mathbf{L}\) and \(\mathbf{W}\) a neighbourhood of \(e\) in \(\mathbf{G}\) such that \(0 = \exp_0 | \mathbf{V}\) is an analytic isomorphism of \(\mathbf{V}\) onto \(\mathbf{W}\) . Let \(\psi = \theta^{-1} \circ (\phi | \phi^{-1}(\mathbf{W}))\) . Then \(\mathrm{T}_0(\psi) = \mathrm{Id}_{\mathrm{L}}\) , whence \(\psi = \mathrm{Id}_{\mathrm{L}}\) using the equality \(\psi\left(\frac{1}{n}x\right) = \frac{1}{n}\psi(x)\) for \(x\) sufficiently close to 0 and \(n \in \mathbf{Z} - \{0\}\) .
\item
  Let \(\mathfrak{a}_1\) be the commutative Lie algebra \(\mathbf{R}^4\) , identified with \(\mathbf{C}^2\) . Let \(\mathfrak{a}_2\) be the commutative Lie algebra \(\mathbf{R}\) . Let \(\phi\) be the homomorphism of \(\mathfrak{a}_2\) into \(\mathrm{Der}(\mathfrak{a}_1)\) which maps 1 to the derivation \((z_1, z_2) \mapsto (iz_1, i\sqrt{2}z_2)\) . Let \(\mathfrak{a}\) be the semi-direct product of \(\mathfrak{a}_2\) by \(\mathfrak{a}_1\) corresponding to \(\phi\) .
\end{enumerate}

\begin{enumerate}
\def\labelenumi{(\alph{enumi})}
\tightlist
\item
  Show that (Int a) \textbar{} \(a_{1}\) contains, for all \(\phi \in R\) , the automorphism
\end{enumerate}

\[
\left(z _ {1}, z _ {2}\right) \mapsto \left(e ^ {i \phi} z _ {1}, e ^ {i \sqrt {2} \phi} z _ {2}\right).
\]

\begin{enumerate}
\def\labelenumi{(\alph{enumi})}
\setcounter{enumi}{1}
\item
  Show that the closure G of (Int a) \textbar{} a₁ in GL(a₁) contains, for all (φ, φ') ∈ R², the automorphism (z₁, z₂) ↦ (eⁱΦz₁, eⁱΦ'z₂).
\item
  Show that \(\mathbf{L}(\mathbf{G})\) contains the endomorphism \(u: (z_1, z_2) \mapsto (z_1, 0)\) of \(\mathfrak{a}_1\) and that there exists no \(x \in \mathfrak{a}\) such that \(u = (\operatorname{ad} x) \mid \mathfrak{a}_1\) .
\item
  Deduce that \(\operatorname{Int}(\mathfrak{a})\) is not closed in \(\operatorname{Aut}(\mathfrak{a})\) .
\end{enumerate}

\begin{enumerate}
\def\labelenumi{\arabic{enumi}.}
\setcounter{enumi}{14}
\tightlist
\item
  Show that the finite-dimensional simply connected real Lie group of §3, Exercise 7 (a) contains non-simply connected Lie subgroups (it contains in fact Lie groups isomorphic to U).
\end{enumerate}

  ¶ 16. (a) Let g be a complex Lie algebra. Let \(\bar{g}\) be the complex Lie algebra derived from g using the automorphism \(\lambda \mapsto \bar{\lambda}\) of C. Let \(g_0\) be the real Lie algebra
derived from g by restriction of scalars. Let \(g' \supset g_0\) be the complex Lie algebra \(g_0 \otimes_{R} C\) . For all \(x \in g\) , we write

\[
f (x) = \frac {1}{2} (x \otimes 1 - (i x) \otimes i) \in \mathfrak {g} ^ {\prime}, \quad g (x) = \frac {1}{2} (x \otimes 1 + (i x) \otimes i) \in \mathfrak {g} ^ {\prime}.
\]

Show that \(f(\text{resp. } g)\) is an isomorphism of \(\mathfrak{g}\) (resp. \(\overline{\mathfrak{g}}\) ) onto an ideal \(\mathfrak{m}\) (resp. \(\mathfrak{n}\) ) of \(\mathfrak{g}'\) and that the ideals \(\mathfrak{m}, \mathfrak{n}\) are supplementary in \(\mathfrak{g}'\) . This defines projectors of \(\mathfrak{g}'\) onto \(\mathfrak{m}, \mathfrak{n}\) . Show that, for all \(x \in \mathfrak{g}, f(x) = p(x), g(x) = q(x)\) .

\begin{enumerate}
\def\labelenumi{(\alph{enumi})}
\setcounter{enumi}{1}
\item
  Suppose that g is finite-dimensional. Let G be the simply connected complex Lie group with Lie algebra g. Let \(\overline{G}\) be the conjugate complex Lie group and \(G_{0}\) the underlying real Lie group. Then \(\mathrm{L}(\overline{\mathbf{G}})=\overline{\mathfrak{g}}\) and \(\mathrm{L}(G_{0})=g_{0}\) . Let M (resp. N) be the simply connected complex Lie group with Lie algebra m (resp. n). Then f defines an isomorphism \(\phi\) of G onto M, g defines an isomorphism \(\gamma\) of \(\overline{G}\) onto N and the simply connected complex Lie group \(G'\) with with Lie algebra \(g'\) is identified with M × N. Show that the real integral subgroup of \(G'\) with Lie algebra \(g_{0}\) is closed and simply connected and is identified with \(G_{0}\) . The restriction to G ⊆ M × N of \(pr_{1}\) (resp. \(pr_{2}\) ) is the isomorphism \(\phi\) (resp. \(\gamma\) ) of G (resp. \(\overline{G}\) ) onto M (resp. N).
\item
  Deduce from (b) and Remark 2 of no. 10 that \(\mathbf{G}'\) , with the canonical injection of \(\mathbf{G}_0\) into \(\mathbf{G}'\) , is the universal complexification of \(\mathbf{G}_0\) .
\item
  Let H be a connected complex Lie group with Lie algebra g, \(\bar{H}\) the conjugate complex Lie group and \(H_{0}\) the underlying real Lie group, such that G is the universal covering space of H. Let K be the (discrete) kernel of the canonical morphism of G onto H. Show that K is a central subgroup of \(G'\) . The canonical injection of \(G_{0}\) into \(G'\) therefore defines an injection i of \(H_{0}\) into \(G'/N\) . Show that \((G'/N, i)\) is the universal complexification of \(H_{0}\) . (Note that this ordered pair has the universal property of Proposition 20.)
\item
  We write \(G^{\prime}/N = (H_{0})_{\mathbf{c}}\) . Deduce from (d) a canonical morphism \(\psi\) of \((\mathrm{H}_{0})_{\mathbf{c}}\) onto \(H \times \bar{H}\) , which is a covering mapping. Show, by the example \(H = C^{*}\) , that \(\psi\) is not in general an isomorphism.
\end{enumerate}

\begin{enumerate}
\def\labelenumi{\arabic{enumi}.}
\setcounter{enumi}{16}
\tightlist
\item
  \begin{enumerate}
  \def\labelenumii{(\alph{enumii})}
  \tightlist
  \item
    Let G, \(\overline{G}\) , \(G_{0}\) be as in Exercise 16 (b). Let V be a finite-dimensional complex vector space and \(\rho\) an irreducible linear representation of \(G_{0}\) on V. Show that there exist finite-dimensional complex vector spaces X, Y, an analytic irreducible linear representation \(\sigma\) (resp. \(\tau\) ) of G (resp. \(\overline{G}\) ) on X (resp. Y) and an isomorphism of V onto X \(\otimes\) Y which transforms \(\rho\) into \(\sigma \otimes \tau\) . (Use Exercise 16 (a), Proposition 2 of Chapter I, § 2, and the Corollary to Proposition 8 of Algebra, Chapter VIII, § 7.)
  \end{enumerate}
\end{enumerate}

\begin{enumerate}
\def\labelenumi{(\alph{enumi})}
\setcounter{enumi}{1}
\tightlist
\item
  Show that the conclusion of (a) is not necessarily true if we omit the hypothesis that G is simply connected. (Consider the complex Lie group \(\mathbf{C}^*\) .)
\end{enumerate}

\begin{enumerate}
\def\labelenumi{\arabic{enumi}.}
\setcounter{enumi}{17}
\tightlist
\item
  Let \(\Lambda\) be a finite-dimensional connected commutative complex Lie group with Lie algebra \(a\) ; let \(\Lambda\) be the kernel of \(\exp_{A}\) , so that \(\Lambda\) is identified with \(a / \Lambda\) .
\end{enumerate}

\begin{enumerate}
\def\labelenumi{(\alph{enumi})}
\tightlist
\item
  The following conditions are equivalent:
\end{enumerate}

(al) The canonical mapping \(\mathbf{C} \otimes_{\mathbf{Z}} \Lambda \to a\) is injective.

(a2) A is isomorphic to a Lie subgroup of some \((\mathbf{C}^{*})^{n}\) .

(a3) A is isomorphic to a group \((\mathbf{C}^{*})^{p}\times \mathbf{C}^{q}\) .

(a4) A has a finite-dimensional faithful complex linear representation.

(a5) A has a finite-dimensional semi-simple faithful complex linear representation with closed image.

\begin{enumerate}
\def\labelenumi{(\alph{enumi})}
\setcounter{enumi}{1}
\tightlist
\item
  The following conditions are equivalent:
\end{enumerate}

(b1) The canonical mapping \(\mathbf{C} \otimes_{\mathbf{Z}} \Lambda \to \mathfrak{a}\) is surjective.

(b2) A is isomorphic to a quotient of a group \((\mathbf{C}^{*})^{n}\) .

(b3) No direct factor of A is isomorphic to C.

(b4) Every complex linear representation of A is semi-simple.

\begin{enumerate}
\def\labelenumi{(\alph{enumi})}
\setcounter{enumi}{2}
\tightlist
\item
  The following conditions are equivalent:
\end{enumerate}

\begin{enumerate}
\def\labelenumi{(\roman{enumi})}
\setcounter{enumi}{149}
\tightlist
\item
  The canonical mapping \(\mathbf{C} \otimes_{\mathbf{Z}} \Lambda \to \mathfrak{a}\) is bijective.
\end{enumerate}

(c2) A is isomorphic to a group \((\mathbf{C}^{*})^{n}\) .

\begin{enumerate}
\def\labelenumi{(\alph{enumi})}
\setcounter{enumi}{3}
\tightlist
\item
  Let \(\mathbf{F}\) be a finite subgroup of \(\mathbf{A}\) and let \(\mathbf{A}' = \mathbf{A} / \mathbf{F}\) . Show that \(\mathbf{A}\) satisfies conditions \((a_i)\) (resp. \((b_i), (c_i)\) ) if and only if \(\mathbf{A}'\) does.
\end{enumerate}

\begin{enumerate}
\def\labelenumi{\arabic{enumi}.}
\setcounter{enumi}{18}
\tightlist
\item
  Let \(G\) be a real Lie group. Show that there exists a neighbourhood \(V\) of 0 in \(L(G)\) such that, for \(x, y\) in \(V\) :
\end{enumerate}

\[
\exp (t (x + y)) = \lim _ {n \rightarrow + \infty} \left(\exp \frac {t x}{n}. \exp \frac {t y}{n}\right) ^ {n}
\]

\[
\exp (t ^ {2} [ x, y ]) = \lim _ {n \rightarrow + \infty} \left(\exp \frac {t x}{n}. \exp \frac {t y}{n}. \exp \frac {- t x}{n}. \exp \frac {- t y}{n}\right) ^ {n ^ {2}}
\]

uniformly for \(t \in (0, 1)\) .

\begin{enumerate}
\def\labelenumi{\arabic{enumi}.}
\setcounter{enumi}{19}
\item
  Let G be a Lie group, H a Lie subgroup of G and A an integral subgroup of G such that \(\mathrm{L}(\mathrm{H}) \cap \mathrm{L}(\mathrm{A}) = \{0\}\) . Show that \(H \cap A\) is discrete in the Lie group A.
\item
  Let \(G\) be a real Lie group and \(H\) a normal 1-dimensional integral subgroup.
\end{enumerate}

\begin{enumerate}
\def\labelenumi{(\alph{enumi})}
\item
  If H is not closed, \(\bar{H}\) is compact (Spectral Theories, Chapter II, §2, Lemma 1), hence isomorphic to \(T^{n}\) and hence central in G if G is connected (use General Topology, Chapter VII, §2, Proposition 5).
\item
  Suppose that H is closed. Let a be an element of G and C(a) the set of elements of G which commute with a. Suppose that H ⊕ C(a).
\end{enumerate}

If H is isomorphic to R, then \(\mathrm{C}(a) \cap \mathrm{H} = \{e\}\) . (Consider the automorphism \(\alpha: h \mapsto a^{-1}ha\) of H.) If H is isomorphic to T, \(\mathrm{C}(a) \cap \mathrm{H}\) has two elements. (Consider \(\alpha\) again and use General Topology, Chapter VII, §2, Proposition 6.) The second situation is impossible if G is connected (use \((a)\) ).

\begin{enumerate}
\def\labelenumi{(\alph{enumi})}
\setcounter{enumi}{2}
\tightlist
\item
  With the hypotheses of (b), suppose further that G/H is commutative.
\end{enumerate}

Then \(\mathrm{G} = \mathrm{C}(a)\) . H. (Observe that the mapping \(h \mapsto h^{-1}\alpha(h)\) of H into H is surjective.)

\begin{enumerate}
\def\labelenumi{\arabic{enumi}.}
\setcounter{enumi}{21}
\tightlist
\item
  Let G be a real Lie group, H a normal 1-dimensional integral subgroup and A a closed subgroup of G such that AH is not closed in G. Then H is central in the identity component of AH. (Reduce it to the case where G = AH and G is connected. Then the set B of elements of A which commute with H is a normal subgroup of G. Passing to the quotient by B, reduce it to the case B = \(\{e\}\) . As every commutator in G commutes with H, A is then commutative. By assuming that H is closed and using Exercise 21 (c), show that AH would be closed, whence a contradiction. Hence H is not closed and Exercise 21 (a) can be used.)
\end{enumerate}

¶ 23. Let G be a real Lie group and \(c = (\mathrm{U}, \phi, \mathrm{E})\) a chart on the manifold G centred on the identity element \(e\) . If \(x \in \mathrm{U}\) , let \(|x|_c\) denote the norm of \(\phi(x)\) in the Banach space E.

\begin{enumerate}
\def\labelenumi{(\alph{enumi})}
\tightlist
\item
  Show that, if \(c' = (\mathbf{U}', \phi', \mathrm{E}')\) is another chart centred on \(e\) , there exist constants \(\lambda > 0\) and \(\mu > 0\) such that
\end{enumerate}

\[
\mu | x | _ {c} \leqslant | x | _ {c ^ {\prime}} \leqslant \lambda | x | _ {c}
\]

for all \(x \in G\) sufficiently close to e.

\begin{enumerate}
\def\labelenumi{(\alph{enumi})}
\setcounter{enumi}{1}
\tightlist
\item
  Show that, for all \(\rho > 0\) , there exists a neighbourhood \(\mathbf{U}_{\rho}\) of \(e\) contained in \(\mathbf{U}\) such that, for \(x, y\) in \(\mathbf{U}_{\rho}\) , \((x, y) \in \mathbf{U}_{\rho}\) and
\end{enumerate}

\[
\left| (x, y) \right| _ {c} \leqslant \rho . \inf (\left| x \right| _ {c}, \left| y \right| _ {c}).
\]

\begin{enumerate}
\def\labelenumi{(\alph{enumi})}
\setcounter{enumi}{2}
\tightlist
\item
  Suppose that \(\mathbf{G}\) is finite-dimensional. Let \(\Gamma\) be a discrete subgroup of \(\mathbf{G}\) . Apply (b) with \(0 < \rho < 1\) and choose \(\mathrm{U}_{\rho}\) to be relatively compact. The set \(\mathrm{U}_{\rho} \cap \Gamma\) is finite. Let
\end{enumerate}

\[
\{e = \gamma_ {0}, \gamma_ {1}, \dots , \gamma_ {m} \}
\]

be its elements, numbered so that

\[
\left| \gamma_ {0} \right| _ {c} \leqslant \left| \gamma_ {1} \right| _ {c} \leqslant \dots \leqslant \left| \gamma_ {m} \right| _ {c}.
\]

Show that, if \(i \leqslant m\) and \(j \leqslant m\) , the commutator \((\gamma_{i}, \gamma_{j})\) is equal to one of the \(\gamma_{k}\) with \(k < \inf(i,j)\) . Deduce that the subgroup generated by \(\mathrm{U}_{\rho} \cap \Gamma\) is nilpotent of class \(\leqslant m\) .

Deduce from the above the existence of a neighbourhood \(\mathbf{V}\) of \(e\) such that, for every discrete subgroup \(\Gamma\) of \(\mathbf{G}\) , there exists a nilpotent integral subgroup \(\mathbf{N}\) of \(\mathbf{G}\) containing \(\mathbf{V} \cap \Gamma\) .

\begin{enumerate}
\def\labelenumi{(\alph{enumi})}
\setcounter{enumi}{3}
\tightlist
\item
  Suppose that G is connected and finite-dimensional and contains an increasing sequence of discrete subgroups \(D_{n}\) whose union is dense in G. Then G is nilpotent. (Using (c), prove the existence of a central one-parameter subgroup H which meets \(D_{n}\) for sufficiently large n at a point distinct from e. Argue by induction on dim G, distinguishing two cases according to whether H is closed or relatively compact (Exercise 21 (a)).
\end{enumerate}

\begin{enumerate}
\def\labelenumi{\arabic{enumi}.}
\setcounter{enumi}{23}
\tightlist
\item
  Let \(G, G'\) be finite-dimensional real Lie groups, \(f\) a surjective morphism of \(G\) into \(G'\) , \(N\) its kernel, \(H\) an integral subgroup of \(G\) and \(H' = f(H)\) .
\end{enumerate}

\begin{enumerate}
\def\labelenumi{(\alph{enumi})}
\item
  Suppose that N is finite. For \(H'\) to be closed, it is necessary and sufficient that H be closed.
\item
  Suppose that N is compact. If H is closed, \(H'\) is closed.
\end{enumerate}

\begin{enumerate}
\def\labelenumi{\arabic{enumi}.}
\setcounter{enumi}{24}
\tightlist
\item
  Let \(G\) be a finite-dimensional real or complex Lie group. Let \(S \subset L(G)\) . Suppose that \(L(G)\) is the Lie algebra generated by \(S\) and that \(S\) is stable under homotheties.
\end{enumerate}

\begin{enumerate}
\def\labelenumi{(\alph{enumi})}
\item
  Let H be the subgroup of G generated by exp S. Show that H is open in G. (Let A be the set of \(x \in L(G)\) such that \(\exp(Kx) \subset H\) and B the vector subspace of \(L(G)\) generated by A. Show that {[}A, S{]} \(\subset A\) and then that {[}A, A{]} \(\subset B\) . Deduce that B is a subalgebra of \(L(G)\) and then use Proposition 3 of § 4.)
\item
  Suppose further that
\end{enumerate}

\[
(x \in S \text {   and   } y \in S) \Rightarrow ((\operatorname{Ad} \exp x) (y) \in S).
\]

Then the vector subspace \(\mathbf{V}\) of \(\mathbf{L}(\mathbf{G})\) generated by \(\mathbf{S}\) is equal to \(\mathbf{L}(\mathbf{G})\) . (Using \((a)\) , show that \([\mathbf{L}(\mathbf{G}), \mathbf{V}] \subset \mathbf{V}\) .)

\begin{enumerate}
\def\labelenumi{\arabic{enumi}.}
\setcounter{enumi}{25}
\tightlist
\item
  Let G be an n-dimensional connected compact complex Lie group, X a finite-dimensional connected complex analytic manifold and \((g, x) \mapsto gx\) a law of analytic left operation of G on X. For all \(g \in G\) , let \(\rho(g)\) be the mapping \(x \mapsto gx\) of X into X. Suppose, that, for all \(g \in G\) distinct from \(e, \rho(g) \neq Id_x\) . Then every orbit of G in X is an n-dimensional closed submanifold of X. (Let \(x \in X\) , H be the stabilizer of x and \(H'\) the identity component of H. For all \(g \in H'\) , let \(u(g)\) be the tangent mapping at x to \(\rho(g)\) . Arguing as in Proposition 6 of no. 3, show that \(u(g)\) is the identity for all \(g \in H'\) . Using § 1, Exercise 4 and the connectedness of X, deduce that \(H' = \{e\}\) .)
\end{enumerate}

¶ 27. Let G be a compact real Lie group, M a compact manifold of class C² and J an open interval of R containing 0. Let \((s, \langle x, \xi \rangle) \mapsto (m_{\xi}(s, x), \xi)\) be a mapping of class C² of G × (M × J) into M × J under which G operates on M × J on the left.

\begin{enumerate}
\def\labelenumi{(\alph{enumi})}
\item
  Let X be a vector field of class \(C^{1}\) on \(M \times J\) such that, for all \((x, \xi) \in M \times J\) , the second projection of \(\mathbf{X}_{(x, \xi)}\) is the tangent vector l to J. Translating X by G and integrating over G, deduce the existence of a field \(X'\) with the same properties as X, which is moreover invariant under G.
\item
  Show that there exists a diffeomorphism \((x, \xi) \mapsto (h_{\xi}(x), \xi)\) of \(\mathbf{M} \times \mathbf{J}\) onto itself such that
\item
  for all \(\xi \in \mathbf{J}\) , \(h_{\xi}\) is a diffeomorphism of \(\mathbf{M}\) onto \(\mathbf{M}\) ;
\end{enumerate}

\begin{enumerate}
\def\labelenumi{(\roman{enumi})}
\setcounter{enumi}{1}
\tightlist
\item
  \(m_{\xi}(s,x) = h_{\xi}(m_0(s,h_{\xi}^{-1}(x)))\) for \(s\in \mathbf{G}, x\in \mathbf{M},\xi \in \mathbf{J}\) .
\end{enumerate}

(Use (a) and Theorem 5 of no. 8.)

\begin{enumerate}
\def\labelenumi{\arabic{enumi}.}
\setcounter{enumi}{27}
\item
  Give an example of a finite-dimensional real Lie group G and two integral subgroups A, B of G such that \(\mathrm{L}(\mathrm{G})=\mathrm{L}(\mathrm{A})+\mathrm{L}(\mathrm{B})\) but AB is not a subgroup of G (cf.~Integration, Chapter VII, § 3, no. 3, Proposition 6).
\item
  Let G be a finite-dimensional complex connected Lie group, \(G_{0}\) the underlying real Lie group and H an integral subgroup of \(G_{0}\) . Show that there exists a smallest integral subgroup \(H^{*}\) of G containing H. Give an example where H is closed in \(G_{0}\) but \(H^{*}\) is not closed in G (take \(G = C^{2}/Z^{2}\) ).
\item
  Let \(G\) be a compact connected complex Lie group and hence of the form \(\mathbf{C}^n / D\) , where \(D\) is a discrete subgroup of \(\mathbf{C}^n\) of rank \(2n\) .
\end{enumerate}

\begin{enumerate}
\def\labelenumi{(\alph{enumi})}
\item
  Show that every holomorphic differential 1-form \(\omega\) on G is invariant. (Let \(\pi: \mathbf{C}^n \to \mathbf{G}\) be the canonical morphism. Let \(\zeta_1, \ldots, \zeta_n\) be the coordinate functions on \(\mathbf{C}^n\) . Then \(\pi^*(\omega)\) is of the form \(\sum_j a_j d\zeta_{j}\) , where the \(a_j\) are holomorphic functions on \(\mathbf{C}^n\) which are invariant under D and hence constant.)
\item
  Let \(G' = C^{n}/D'\) , where \(D'\) is a discrete subgroup of \(C^{n}\) of rank 2n. Show that every analytic manifold isomorphism \(u: G \to G'\) is of the form \(s \mapsto v(s) + a'\) , where v is a Lie group isomorphism and \(a' \in G'\) . (Let \(\pi': C^{n} \to G'\) be the canonical morphism. There exists an isomorphism of analytic manifolds \(\tilde{u}: C^{n} \to C^{n}\) such that \(u \circ \pi' \circ \tilde{u}\) . For every holomorphic differential 1-form \(\omega'\) on \(G'\) , \(u^*(\pi^*(\omega'))\) is a 1-form on \(C^{n}\) which is invariant under translations by (a). Deduce that \(\tilde{u}\) is an affine mapping.)
\end{enumerate}

\exercisesfor{7}

\begin{enumerate}
\def\labelenumi{\arabic{enumi}.}
\item
  Let \(\mathrm{G}\) be a Lie group and \(\phi: \mathrm{U} \to \mathrm{G}\) an exponential mapping such that \(\mathbf{Z}\mathbf{U} \subset \mathbf{U}\) and \(\phi(rx) = \phi(x)^r\) for all \(x \in \mathbf{U}\) and all \(r \in \mathbf{Z}\) . If \(p > 0\) , \(\phi\) is an analytic isomorphism of \(\mathrm{U}\) onto \(\phi(\mathrm{U})\) . (If \(\phi(x) = \phi(y)\) , then \(\phi(p^n x) = \phi(p^n y)\) for all \(n \in \mathbf{N}\) and hence \(x = y\) . Let \(\mathrm{W}\) be an open neighbourhood of \(0\) in \(\mathrm{L}(\mathrm{G})\) such that \(\phi^{-1}\) is analytic on \(\phi(\mathrm{W})\) . For all \(s \in \phi(\mathrm{U})\) , there exist \(n \in \mathbf{N}\) and a neighbourhood \(\mathrm{V}\) of \(s\) in \(\phi(\mathrm{U})\) such that \(t \in \mathrm{V} \Rightarrow t^{p^n} \in \phi(\mathrm{W})\) .)
\item
  Let G be a Lie group, \(\phi: \mathrm{U} \to \mathrm{G}\) an exponential mapping with the properties of Proposition 3, V a neighbourhood of 0 in U such that \(\mathbf{ZV} \subset \mathrm{V}\) and \(\psi: \mathrm{V} \to \mathrm{G}\) a tangent mapping at 0 to \(\phi\) such that \(\psi(nx) = \psi(x)^n\) for all \(n \in \mathrm{V}\) and all \(n \in \mathbf{Z}\) . If \(p > 0\) , then \(\psi = \phi \mid \mathrm{V}\) . (Give L(G) a norm. Let \(x \in \mathrm{V}\) . Then \(p^n x\) tends to 0 as \(n\) tends to \(+\infty\) , hence there exists \(\alpha_n > 0\) such that \(\alpha_n\) tends to 0 and \(\|(\phi^{-1} \circ \psi)(p^n x) - p^n x \| \leqslant \alpha_n \|p^n x\|\) . But \(\psi(p^n x) = \psi(x)p^n\) , whence \(\|(\phi^{-1} \circ \psi)(x) - x \| \leqslant \alpha_n \|x\|\) .)
\item
  Let U be the set of invertible elements of A and \(U' = 1 + m \subset U\) .
\end{enumerate}

\begin{enumerate}
\def\labelenumi{(\alph{enumi})}
\item
  Show that \(\mathbf{U}' \subset \mathbf{U}_f\) . (If \(x = 1 + y\) with \(y \in \mathfrak{m}\) , then \(x^{p^n}\) tends to 1 as \(n\) tends to \(+\infty\) by the binomial theorem.)
\item
  Show that \(U_{f}\) is the set of elements of U whose image in A/m is a root of unity. (Use (a).) Hence recover the fact that \(U = U_{f}\) when K is locally compact (A/m is then finite).
\end{enumerate}

\begin{enumerate}
\def\labelenumi{\arabic{enumi}.}
\setcounter{enumi}{3}
\tightlist
\item
  Let \(n \in \mathbf{N}^*\) , \(p\) be a prime number and \(G\) the set of matrices belonging to \(\mathbf{GL}(n, \mathbf{Z}_p)\) all of whose elements are congruent to 1 modulo \(p\) if \(p \neq 2\) (resp. modulo 4 if \(p = 2\) ). Then \(G\) is an open subgroup of \(\mathbf{GL}(n, \mathbf{Z}_p)\) .
\end{enumerate}

\begin{enumerate}
\def\labelenumi{(\alph{enumi})}
\item
  Show that G has no element of finite order \(\neq 1\) .
\item
  Show that every finite subgroup of \(\mathbf{GL}(n, \mathbf{Z}_p)\) is isomorphic to a subgroup of \(\mathbf{GL}(n, \mathbf{Z} / p\mathbf{Z})\) if \(p \neq 2\) (resp. \(\mathbf{GL}(n, \mathbf{Z} / 4\mathbf{Z})\) if \(p = 2\) ).
\end{enumerate}

¶ 5. (a) Let \(\Gamma\) be a compact subgroup of \(G = \mathbf{GL}(n, \mathbf{Q}_p)\) . Show that there exists a conjugate of \(\Gamma\) which is contained in \(\mathbf{GL}(n, \mathbf{Z}_p)\) . (If \(T\) is a lattice of \(\mathbf{Q}_p^n\) with respect to \(\mathbf{Z}_p\) (Commutative Algebra, Chapter VII, § 4, Definition 1), show that the stabilizer of \(T\) in \(\Gamma\) is open in \(\Gamma\) and hence of finite index and \(\sum_{\gamma \in I} \gamma T\) is a lattice which is stable under \(\Gamma\) .)

\begin{enumerate}
\def\labelenumi{(\alph{enumi})}
\setcounter{enumi}{1}
\item
  Deduce that \(\mathbf{G}_f\) is the union of the conjugates of \(\mathbf{GL}(n,\mathbf{Z}_p)\) .
\item
  Deduce that every finite subgroup \(\Phi\) of \(\mathbf{GL}(n, \mathbf{Q}_p)\) has order a divisor of \(a_n(p)\) , where \(a_n(p)\) is defined by
\end{enumerate}

\[
\begin{array}{l l} a _ {n} (p) = (p ^ {n} - 1) (p ^ {n} - p) \ldots (p ^ {n} - p ^ {n - 1}) & \text {if} p \neq 2 \\ a _ {n} (p) = 2 ^ {n ^ {2}} (2 ^ {n} - 1) (2 ^ {n} - 2) \ldots (2 ^ {n} - 2 ^ {n - 1}) & \text {if} p = 2. \end{array}
\]

(Using (a), it can be assumed that \(\Phi \subset \mathbf{GL}(n, \mathbf{Z}_p)\) . Then use Exercise 4.)

\begin{enumerate}
\def\labelenumi{(\alph{enumi})}
\setcounter{enumi}{3}
\tightlist
\item
  Show that every finite subgroup of \(\mathbf{SL}(n, \mathbf{Q}_p)\) has order a divisor of \(s_n(p)\) , where \(s_n(p) = \frac{a_n(p)}{p - 1}\) for \(p \neq 2\) and \(s_n(2) = \frac{a_n(2)}{2}\) .
\end{enumerate}

¶ 6. In this exercise, l denotes a prime number \(\neq2\) . If a is an integer \(\neq0\) , let \(v_{l}(a)\) denote the l-adic valuation of a, i.e.~the largest integer e such that \(a\equiv0\pmod{l^{e}}\) .

\begin{enumerate}
\def\labelenumi{(\alph{enumi})}
\tightlist
\item
  Let \(m\) be an integer \(\geqslant 1\) . We write
\end{enumerate}

\[
\begin{array}{l l} \varepsilon (l, m) = 0 & \text { if } m \neq 0 (\mathrm{mod} (l - 1)) \\ \varepsilon (l, m) = v _ {l} \left(\frac {m}{l - 1}\right) + 1 & \text { if } m \equiv 0 (\mathrm{mod} (l - 1)). \end{array}
\]

Show that, if \(x\) is an integer prime to \(l\) , then

\[
v _ {l} (x ^ {m} - 1) \geqslant \varepsilon (l, m)
\]

and that there is equality if the image of \(x\) in the cyclic group \((\mathbf{Z} / l^2\mathbf{Z})^*\) is a generator of this group.

\begin{enumerate}
\def\labelenumi{(\alph{enumi})}
\setcounter{enumi}{1}
\tightlist
\item
  Let \(n\) be an integer \(\geqslant 1\) . We write
\end{enumerate}

\[
r (l, n) = \sum_ {m = 1} ^ {n} \varepsilon (l, m).
\]

Show that

\[
r (l, n) = \left[ \frac {n}{l - 1} \right] + \left[ \frac {n}{l (l - 1)} \right] + \left[ \frac {n}{l ^ {2} (l - 1)} \right] + \dots
\]

where the symbol \([\alpha]\) denotes the integral part of the real number \(\alpha\) . Show that, if x is an integer prime to l, then

\[
v _ {l} ((x ^ {n} - 1) (x ^ {n - 1} - 1) \dots (x - 1)) \geqslant r (l, n)
\]

and that there is equality if the image of \(x\) in \((\mathbf{Z} / l^2\mathbf{Z})^*\) is a generator of this group.

\begin{enumerate}
\def\labelenumi{(\alph{enumi})}
\setcounter{enumi}{2}
\item
  In the notation of Exercise 5 (c), show that \(l^{r(l,n)}\) is the greatest power of \(l\) which divides all the \(a_{n}(p)\) , for \(p\) prime \(\neq 1\) (or for all sufficiently large \(p\) , which amounts to the same). (Use (b) applied to \(x = p\) ; then choose \(p\) such that its image in \((\mathbf{Z}/l^2\mathbf{Z})^*\) is a generator of this group, which is possible by the arithmetic progression theorem.†)
\item
  Let \(\bar{\Gamma}\) be a finite subgroup of \(\mathbf{GL}(n, \mathbf{Q})\) and let \(l^e\) be the greatest power of \(l\) dividing the order of \(\Gamma\) . Show that \(e \leqslant r(l, n)\) . (Use Exercise 5 to prove that \(l^e\) divides all the \(a_n(p)\) and then apply (c) above.)
\item
  Conversely, show that there exists a finite \(l\) -subgroup \(\Gamma_{i,n}\) of \(\mathbf{GL}(n,\mathbf{Q})\) whose order is \(l^{d,l,n}\) . (Reduce it to the case where \(n\) is of the form \(l^a (l - 1)\) , with \(a\geqslant 0\) . Decompose \(\mathbf{Q}^{\mathrm{n}}\) as a sum of \(l^a\) copies of \(\mathbf{Q}^{j - 1}\) and use this decomposition to give an operation on \(\mathbf{Q}^{\mathrm{n}}\) of the semi-direct product \(\mathrm{H}_{l,n}\) of the symmetric group \(\mathfrak{S}_{la}\) and the group \((\mathbf{Z} / l\mathbf{Z})^{a}\) . Take \(\Gamma_{l,n}\) to be a Sylow \(l\) -group of \(\mathrm{H}_{l,n}\) .)
\end{enumerate}

Show that, if \(n\) is even, \(\Gamma_{l,n}\) is contained in a conjugate of the symplectic group \(\mathbf{Sp}(n,\mathbf{Z})\) .

\begin{enumerate}
\def\labelenumi{(\alph{enumi})}
\setcounter{enumi}{5}
\tightlist
\item
  *Let \(\Gamma\) be a finite subgroup of \(\mathbf{GL}(n, \mathbf{Q})\) which is an \(l\) -group. Show that \(\Gamma\) is conjugate to a subgroup of \(\Gamma_{l,n}\) . (Show first, using a suitable reduction modulo \(p\) , that the reduction of \(\Gamma\) is a finite subgroup of the reduction of a conjugate of \(\Gamma_{l,n}\) ; then use the character of the representation of \(\Gamma\) on \(\mathbf{Q}^{n}\) .) In particular, every subgroup of \(\mathbf{GL}(n, \mathbf{Q})\) of order \(l^{(l,n)}\) is conjugate to \(\Gamma_{l,n*}\) .
\end{enumerate}

¶ 7. In this exercise, let \(v_{2}(a)\) denote the 2-adic valuation of an integer \(a\) .

\begin{enumerate}
\def\labelenumi{(\alph{enumi})}
\item
  Let \(\Gamma\) be a finite subgroup of \(\mathbf{GL}(n,\mathbf{Q})\) . Show that there exists a positive definite quadratic form with coefficients in \(\mathbf{Z}\) , which is invariant under \(\Gamma\) . Deduce (by the same argument as in Exercises 4 and 5) that, for all sufficiently large \(p\) , \(\Gamma\) is isomorphic to a subgroup of an orthogonal group \(\mathbf{O}(n)\) over the field \(\mathbf{F}_p\) (resp. a subgroup of \(\mathbf{SO}(n)\) if \(\Gamma\) is contained in \(\mathbf{SL}(n,\mathbf{Q})\) ).
\item
  Suppose that \(\Gamma\) is contained in \(\mathbf{SL}(n, \mathbf{Q})\) and let \(2^e\) denote the greatest power of 2 dividing the order of \(\Gamma\) . Show that, if \(n\) is odd, \(2^e\) divides all the integers
\end{enumerate}

\[
b _ {n} (p) = \left(p ^ {n - 1} - 1\right) \left(p ^ {n - 3} - 1\right) \dots \left(p ^ {2} - 1\right)
\]

for sufficiently large prime \(p\) . (Use (a) and Exercise 13 of Algebra, Chapter IX, § 6.)

Show that, if \(n\) is even and \(p\) is a sufficiently large prime, \(2^e\) divides the l.c.m. \(b_n(p)\) of the integers

\[
\begin{array}{l} (p ^ {n} - 1) (p ^ {n - 2} - 1) \dots (p ^ {2} - 1) / (p ^ {n / 2} + 1) \\ (p ^ {n} - 1) (p ^ {n - 2} - 1) \dots (p ^ {2} - 1) / (p ^ {n / 2} - 1). \end{array}
\]

(Same method as for \(n\) odd.)

\begin{enumerate}
\def\labelenumi{(\alph{enumi})}
\setcounter{enumi}{2}
\tightlist
\item
  Under the hypotheses of (b), we write
\end{enumerate}

\[
r (2, n) = n + \left[ \frac {n}{2} \right] + \left[ \frac {n}{4} \right] + \dots = n + v _ {2} (n!).
\]

Let A be an integer \(\geqslant 3\) . Show that \(2^{r(2,n)-1}\) is the greatest power of 2 which divides all the \(b_n(p)\) for \(p \geqslant A\) . (Same method† as in Exercise 6; use the existence of a prime number \(p \geqslant A\) such that \(p \equiv 5 \pmod{8}\) .) Deduce the inequality \(e \leqslant r(2,n) - 1\) .

\begin{enumerate}
\def\labelenumi{(\alph{enumi})}
\setcounter{enumi}{3}
\tightlist
\item
  Conversely, let \(C_{n}\) be the subgroup of \(\mathbf{GL}(n,\mathbf{Z})\) generated by the permutation matrices and the diagonal matrices with coefficients \(\pm1\) . The order of \(C_{n}\) is \(2^{n}n!\) and \(v_{2}(2^{n}n!) = r(2,n)\) . If \(\Gamma_{2,n}\) denotes the intersection of a Sylow 2-group of \(C_{n}\) with \(\mathbf{SL}(n,\mathbf{Z})\) , deduce that the order of \(\Gamma_{2,n}\) is \(2^{n(2,n)-1}\) .
\end{enumerate}

\begin{enumerate}
\def\labelenumi{\arabic{enumi}.}
\setcounter{enumi}{7}
\tightlist
\item
  Let \(n\) be an integer \(\geqslant 1\) . We write
\end{enumerate}

\[
\mathrm{M} (n) = \prod_ {l} 1 ^ {r (1, n)},
\]

where the product is taken over all prime numbers \(l\) and the \(r(l, n)\) are defined as in Exercises 6 and 7.

Then \(\mathbf{M}(1) = 2\) , \(\mathbf{M}(2) = 2^3.3 = 24\) , \(\mathbf{M}(3) = 2^4.3 = 48\) ,

\[
\mathrm{M} (4) = 2 ^ {7}. 3 ^ {2}. 5 = 5 7 6 0.
\]

Deduce from Exercises 6 and 7 that the l.c.m. of the orders of the finite subgroups of \(\mathbf{GL}(n,\mathbf{Q})\) (or \(\mathbf{GL}(n,\mathbf{Z})\) , which amounts to the same) is equal to \(\mathrm{M}(n)\) . Consider the same question for \(\mathbf{SL}(n,\mathbf{Q})\) with \(\mathrm{M}(n)\) replaced by \(\frac{1}{2}\mathrm{M}(n)\) .

\begin{enumerate}
\def\labelenumi{\arabic{enumi}.}
\setcounter{enumi}{8}
\item
  Suppose that \(\mathbf{K}\) is locally compact. Let \(\mathrm{G} = \mathbf{GL}(n,\mathbf{K})\) . Let \(\mathbf{G}_1\) be the set of \(g\in \mathbf{G}\) which leaves stable a lattice of \(\mathbf{K}^n\) with respect to A. Let \(\mathbf{G}_2\) be the set of \(g\in \mathbf{G}\) which generate a relatively compact subgroup in G. Let \(\mathbf{G}_3\) be the set of \(g\in \mathbf{G}\) whose eigenvalues in an algebraic closure of \(\mathbf{K}\) are of absolute value 1. Then \(\mathbf{G}_1 = \mathbf{G}_2 = \mathbf{G}_3 = \mathbf{G}_f\) . (Use the argument of Exercise 5 (a).)
\item
  Suppose that K is locally compact. Let G be a standard group of dimension n over K. Let \(\mu\) be a Haar measure on the additive group \(K^{n}\) . Show that \(\mu|G\) is a left and right Haar measure on G. (Use the fact that G is the inverse limit of the \(G(\alpha_{\lambda})\) and Integration, Chapter VII, §1, Proposition 7.)
\end{enumerate}

\exercisesfor{8}

\begin{enumerate}
\def\labelenumi{\arabic{enumi}.}
\item
  The mapping \(z \mapsto \bar{z}\) of \(\mathbf{C}\) into \(\mathbf{C}\) is a non-analytic continuous automorphism of the complex Lie group \(\mathbf{C}\) .
\item
  Let G be the Hilbert space of sequences \((\lambda_{1}, \lambda_{2}, \ldots)\) of real numbers such that \(\sum_{i} \lambda_{i}^{2} < +\infty\) . Consider G as a real Lie group. Let \(G_{n}\) be the set of \((\lambda_{1}, \lambda_{2}, \ldots) \in G\) such that \(\lambda_{m} \in \frac{1}{m} Z\) for \(1 \leqslant m \leqslant n\) . The \(G_{n}\) are closed Lie subgroups of G and hence \(H = \bigcap_{n} G_{n}\) is a closed subgroup of G. But this group is totally disconnected and not discrete and hence is not a Lie subgroup of G.
\end{enumerate}

¶ 3. In \(Q_{p} \times Q_{p}\) , every closed subset can be defined by a family of analytic equations. Deduce that Corollary 2 (ii) of Theorem 2 is no longer true if the hypothesis that I is finite is omitted.

¶ 4. Let G be a finite-dimensional real Lie group and A a subgroup of G. We say that an element x of L(G) is A-accessible if, for every neighbourhood U of e, there exists a continuous mapping \(\alpha\) of {[}0, 1{]} into A such that \(\alpha(0) = e\) and \(\alpha(t) \in \exp(tx)\) . U for \(0 \leqslant t \leqslant 1\) . Let b be the set of A-accessible elements of L(G).

\begin{enumerate}
\def\labelenumi{(\alph{enumi})}
\item
  Show that \(\mathfrak{h}\) is a Lie subalgebra of \(\mathbf{L}(\mathbf{G})\) . (Use Exercise 19 of §6.)
\item
  Let H be an integral subgroup of G such that \(\mathrm{L}(\mathrm{H}) = \mathfrak{h}\) . Show that \(H \subset A\) . (Let I = \(\{-1, 1\}\) and let \(R'\) be given the Euclidean norm. Let \((x_{1}, \ldots, x_{r})\) be a basis of h. Construct continuous mappings \(\alpha_{1}, \ldots, \alpha_{r}\) of I into A and continuous mappings \(f_{1}, \ldots, f_{r}\) from \(I'\) to R, such that, for \(t = (t_{1}, \ldots, t_{r}) \in I'\) ,
\end{enumerate}

\[
\begin{array}{c} \alpha_ {1} (t _ {1}) \ldots \alpha_ {r} (t _ {r}) = \exp (f _ {1} (t) x _ {1}) \ldots \exp (f _ {r} (t) x _ {r}) \\ \| t - (f _ {1} (t), \ldots , f _ {r} (t)) \| \leqslant \frac {1}{2}. \end{array}
\]

Then apply the following theorem: let \(f\) be a continuous mapping of \(\mathbf{I}^r\) into \(\mathbf{R}^r\) such that \(\| f(x) - x \| \leqslant \frac{1}{2}\) for all \(x \in \mathbf{I}^r\) ; then \(f(\mathbf{I}^r)\) contains a neighbourhood of 0 in \(\mathbf{R}^r. \dagger\)

\begin{enumerate}
\def\labelenumi{(\alph{enumi})}
\setcounter{enumi}{2}
\item
  Deduce that \(\mathfrak{h}\) is the subalgebra tangent to A at \(e\) .
\item
  Show that if A is arcwise connected, then A = H.†
\end{enumerate}

¶ 5. Let G be a Hausdorff topological group, H a closed subgroup of G and \(\pi\) the canonical mapping of G onto G/H. Suppose that H is a finite-dimensional real Lie group. There exist a neighbourhood U of \(\pi(e)\) in G/H and a continuous mapping \(\sigma\) of U into G such that \(\pi \circ \sigma = \mathrm{Id}_{\mathrm{U}}\) . (Let \(\rho\) be an analytic linear representation of H on \(\mathbf{GL}(n,\mathbf{R})\) , which is locally a homeomorphism (§ 6, Corollary to Theorem 1). Let \(f\) be a continuous function \(\geqslant 0\) on G, equal to 1 at \(e\) and zero outside a sufficiently small neighbourhood of \(e\) . Let \(ds\) be a left Haar measure of H. For \(x \in G\) , we write

\[
g (x) = \int f (x s) \rho (s) ^ {- 1} d s \in \mathbf {M} _ {n} (\mathbf {R}).
\]

Then \(g(xt) = g(x)\varphi(t)\) for \(x \in G\) and \(t \in H\) . If V is sufficiently small,

\[
g (x) \in \mathbf {G L} (n, \mathbf {R})
\]

for \(x\) sufficiently close to \(e\) . Finally use the fact that the theorem to be proved is true locally for \(\mathbf{GL}(n, \mathbf{R})\) and \(\varphi(\mathbf{H})\) .

\begin{enumerate}
\def\labelenumi{\arabic{enumi}.}
\setcounter{enumi}{5}
\item
  Let G be a group of class \(C^{r}\) (§ 5, Exercise 1) with \(r \geqslant 2\) . There exists on G one and only one real Lie group structure S such that the underlying manifold structure of class \(C^{r}\) of S is the given structure. (The uniqueness of S follows from Corollary 1 to Theorem 1. Let \(L(G)\) be the normable Lie algebra associated with G in Exercise 2 of § 5. There exist a real Lie group germ \(G^{r}\) and an isomorphism h of \(L(G^{\prime})\) onto \(L(G)\) (§ 4, Theorem 3). Verify as in § 4, no. 1 that there exist a symmetric open neighbourhood \(G^{\prime\prime}\) of \(e_{G}\) in \(G^{r}\) and a mapping \(\phi\) of class \(C^{r}\) of \(G^{\prime\prime}\) into G such that \(T_{e}(\phi) = h\) and \(\phi(g_{1}g_{2}) = \phi(g_{1})\phi(g_{2})\) for \(g_{1}, g_{2}\) in \(G^{\prime\prime}\) . By shrinking \(G^{\prime\prime}\) , it can be assumed that \(V = \phi(G^{r})\) is open in G and that \(\phi\) is an isomorphism of class \(C^{r}\) of the manifold \(G^{\prime\prime}\) onto the manifold V. There therefore exists on V a real Lie group germ structure such that the underlying manifold structure of class \(C^{r}\) is the given structure. For all \(g \in G\) , Int g defines as analytic mapping of \(V \cap (g^{-1}Vg)\) onto \((gVg^{-1}) \cap V\) (Theorem 1). By Proposition 18 of § 1, there exists on G a real Lie group structure S inducing the same analytic structure as V on an open neighbourhood of e. By translation, the underlying manifold structure of class \(C^{r}\) of S on G is the given structure.)
\item
  Let \(G\) be a real Lie group and \(H\) a closed subgroup of \(G\) .
\end{enumerate}

\begin{enumerate}
\def\labelenumi{(\alph{enumi})}
\tightlist
\item
  Let \(\mathfrak{h}\) be the set of \(x\in \dot{\mathbf{L}} (\mathbf{G})\) such that \(\exp (tx)\in \dot{\mathbf{H}}\) for all \(t\in \mathbf{R}\) . Then \(\mathfrak{h}\) is a Lie subalgebra of \(\mathrm{L}(\mathrm{G})\) . (Use Proposition 8 of §6.)
\end{enumerate}

\begin{enumerate}
\def\labelenumi{(\alph{enumi})}
\setcounter{enumi}{1}
\tightlist
\item
  Suppose that \(\mathbf{H}\) is locally compact. Show that \(\mathbf{H}\) is a Lie subgroup of \(\mathbf{G}\) . (Show first that \(\mathfrak{h}\) is finite-dimensional by proving the existence in \(\mathfrak{h}\) of a precompact neighbourhood of 0. Then imitate the proof of Theorem 2, choosing \(\mathrm{V}_2\) such that \((\exp \mathrm{V}_2) \cap \mathrm{H}\) is relatively compact.)
\end{enumerate}

\exercisesfor{9}

\begin{enumerate}
\def\labelenumi{\arabic{enumi}.}
\item
  Let \(\mathrm{G} = \mathbf{SO}(3, \mathbf{R})\) . Let \(\Delta_{1}, \Delta_{2}\) be two orthogonal lines in \(R^{3}\) and \(H_{i}\) the subgroup of G consisting of the rotations about \(\Delta_{i} (i = 1, 2)\) . Then \([\mathrm{L}(\mathrm{H}_{1}), \mathrm{L}(\mathrm{H}_{2})]\) is a 1-dimensional Lie subalgebra of G distinct from the Lie subalgebra tangent to \((\mathrm{H}_{1}, \mathrm{H}_{2})\) at e.
\item
  Suppose that K is ultrametric. Let G be a finite-dimensional Lie group, g its Lie algebra and A a finite subset of G. Then \(Z_{G}(A)\) is a Lie subgroup of G with Lie algebra \(\mathfrak{z}_{\mathfrak{g}}(A)\) . (Argue as in Proposition 8.)
\item
  Let G be a connected real or complex Lie group. The centre Z of G is a Lie quasi-subgroup of G and L(Z) is the centre of L(G).
\item
  Suppose that K is ultrametric and p \textgreater{} 0 (in the notation of §7). Let G be a finite-dimensional Lie group, A a group of automorphisms of G and B the corresponding group of automorphisms of L(G). Let \(\mathrm{G}^{\mathbf{A}}\) (resp. \(\mathrm{L}(\mathrm{G})^{\mathbf{B}}\) ) be the set of elements of G (resp. L(G)) which are fixed under A (resp. B). Then \(G^{A}\) is a Lie subgroup of G with Lie algebra \(\mathrm{L}(G)^{\mathbf{B}}\) . (Use the logarithmic mapping.)
\item
  Let \(G\) be a finite-dimensional connected real or complex Lie group. Let \((G_0, G_1, \ldots)\) be the upper central series of (Chapter II, § 4, Exercise 18) and let \((g_0, g_1, \ldots)\) be the upper central series of the Lie algebra \(L(G)\) (Chapter I, § 1, no. 6). Then, for all \(i\) , \(G_i\) is a Lie subgroup of \(G\) such that \(L(G_i) = g_i\) .
\item
  Let \(\Gamma\) be the 3-dimensional simply connected nilpotent real Lie group defined in Exercise 5 (b) of §4. Let \(\alpha \in \mathbf{R}\) be an irrational number. Let \(P\) be the discrete subgroup of \(\Gamma \times \mathbf{R}\) consisting of the \(((0,0,x),\alpha x)\) , where \(x \in \mathbf{Z}\) . Let \(G = (\Gamma \times \mathbf{R}) / P\) . Show that \((G,G)\) is not closed in \(G\) .
\item
  \begin{enumerate}
  \def\labelenumii{(\alph{enumii})}
  \tightlist
  \item
    Let \(G\) be a semi-simple connected real Lie group, \(Z\) its centre and \(\rho\) a continuous linear representation of \(G\) on a finite-dimensional complex vector space. There exists an integer \(p\) such that, for all \(z \in Z\) , \(\rho(z)\) is diagonalizable and all the eigenvalues of \(\rho(z)\) are \(p\) -th roots of unity. (It can be assumed that \(\rho\) is irreducible. Then \(\rho(z)\) is scalar by Schur's Lemma. On the other hand, det \(\rho(g) = 1\) for all \(g \in G\) because \(G = \mathcal{D}G\) .)
  \end{enumerate}
\end{enumerate}

\begin{enumerate}
\def\labelenumi{(\alph{enumi})}
\setcounter{enumi}{1}
\tightlist
\item
  Deduce that, if \(G\) admits an injective finite-dimensional continuous linear representation, \(Z\) is finite.
\end{enumerate}

¶ 8. Let G be a finite-dimensional connected real Lie group, g its Lie algebra, n the largest nilpotent ideal of g and h = {[}g, g{]} + n.

\begin{enumerate}
\def\labelenumi{(\alph{enumi})}
\item
  \(\mathfrak{h}\) is a characteristic ideal of \(\mathfrak{g}\) ; the radical of \(\mathfrak{h}\) is \(n\) ; for all \(x \in \mathfrak{h}\) , \(\operatorname{Tr} \operatorname{ad}_{\mathfrak{h}} x = 0\) ; for every Levi section \(l\) of \(\mathfrak{g}\) , \(\mathfrak{h} = l + n\) .
\item
  Suppose that G is simply connected. Let Z be its centre, H the Lie subgroup of G with Lie algebra b and \(\phi\) the canonical morphism of G onto G/H. Then \(\phi(Z)\) is discrete in G/H. (The group G is the semi-direct product of a semi-simple Lie group S and its radical R. Let \(z \in Z\) . Then \(z = y^{-1}x\) , where \(x \in R\) and y belongs to the centre of S. There exists an integer p such that the eigenvalues of \(Ad_{q}y\) , and hence also of \(Ad_{q}x\) , are p-th roots of unity (Exercise 7). Deduce that, if N is the Lie subgroup of G with Lie algebra n, there exists a neighbourhood U of e in R satisfying U = UN = NU,
\end{enumerate}

\[
\mathrm{Z} \cap (\mathrm{SU}) \subset \mathrm{SN} = \mathrm{H}.
\]

\begin{enumerate}
\def\labelenumi{(\alph{enumi})}
\setcounter{enumi}{2}
\item
  We no longer assume that G is simply connected. Let H be the integral subgroup of G with Lie algebra b. Show that H is a unimodular normal Lie subgroup of G with nilpotent radical, such that G/H is commutative. (Use (a) and (b).)
\item
  Deduce that, if (G, G) is dense in G, the radical of G is nilpotent.
\end{enumerate}

\begin{enumerate}
\def\labelenumi{\arabic{enumi}.}
\setcounter{enumi}{8}
\tightlist
\item
  Let \(G\) be the universal covering of \(\mathbf{SL}(2, \mathbf{R})\) . We identify the kernel of the canonical morphism \(G \to \mathbf{SL}(2, \mathbf{R})\) with \(\mathbf{Z}\) (§ 6, Exercise 2). Let \(a\) be an element of \(\mathbf{T}^n\) whose powers are everywhere dense in \(\mathbf{T}^n\) (General Topology, Chapter VII, §1, Corollary 2 to Proposition 7). Let \(D\) be the discrete subgroup of \(G \times \mathbf{T}^n\) generated by (1, a). Let \(H = (G \times \mathbf{T}^n)/D\) . Then
\end{enumerate}

\[
\mathrm{L} (\mathrm{H}) = \mathfrak {s l} (2, \mathbf {R}) \times \mathbf {R} ^ {n}
\]

and the integral subgroup \(\mathbf{H}'\) of \(\mathbf{H}\) with Lie algebra \(\mathfrak{sl}(2,\mathbf{R})\) is isomorphic to \(\mathbf{G}\) and dense in \(\mathbf{H}\) . \(\mathrm{D}^n\mathrm{H}' = \mathrm{H}'\) for all \(n\geqslant 0\) .

¶ 10. Let G be a finite-dimensional real or complex Lie group. Let

\[
p = \dim [ L (G), L (G) ].
\]

Let \((\mathrm{G},\mathrm{G})\) be given its structure as an integral subgroup of \(\mathbf{G}\) . There exists a neighbourhood \(\mathbf{V}\) of \(e\) in \((\mathrm{G},\mathrm{G})\) such that every element of \(\mathbf{V}\) is the product of \(p\) commutators of elements of \(\mathbf{G}\) . (Let \(x_{1},y_{1},\ldots ,x_{p},y_{p}\) be elements of \(\mathbf{L}(\mathbf{G})\) such that the \([x_1,y_1]\) form a basis of \(\mathbf{L}(\mathbf{G})\) . For \(s,t\) in \(\mathbf{R}\) , we write

\[
\rho_ {i} (s, t) = (\exp s y _ {i}) ^ {- 1} (\exp t x _ {i}) ^ {- 1} (\exp s y _ {i}) (\exp t x _ {i}).
\]

Apply the implicit function theorem to the mapping

\[
\left(s _ {1}, t _ {1}, \dots , s _ {p}, t _ {p}\right) \mapsto \rho_ {1} \left(s _ {1}, t _ {1}\right) \dots \rho_ {p} \left(s _ {p}, t _ {p}\right)
\]

of \(\mathbf{R}^{2p}\) into (G, G).

\begin{enumerate}
\def\labelenumi{\arabic{enumi}.}
\setcounter{enumi}{10}
\item
  Let G be the semi-direct product of B = Z/2Z by K corresponding to the automorphism \(x \mapsto -x\) of the Lie group K. Then L(K) is an ideal of L(G) and L(B) = \{0\}, but (K, B) = K.
\item
  Let \((e_1, e_2, e_3)\) be the canonical basis of \(\mathbf{K}^3\) . Consider the nilpotent Lie algebra structure on \(\mathbf{K}^3\) such that \([e_1, e_2] = e_3\) , \([e_1, e_3] = [e_2, e_3] = 0\) . With the associated group law on \(\mathbf{K}^3\) (no. 5),
\end{enumerate}

\[
(x, y, z) (x ^ {\prime}, y ^ {\prime}, z ^ {\prime}) = (x + x ^ {\prime}, y + y ^ {\prime}, z + z ^ {\prime} + \frac {1}{2} (x y ^ {\prime} - y x ^ {\prime})).
\]

The mapping

\[
(\lambda , \mu , \nu) \mapsto (\exp \lambda e _ {1}) (\exp \mu e _ {2}) (\exp \nu (e _ {1} + e _ {3}))
\]

of \(K^{3}\) into G is neither surjective (show that \((0,1,1)\) is not in the image) nor injective (show that \((0,1,0)\) and \((1,1,-1)\) have the same image).

\begin{enumerate}
\def\labelenumi{\arabic{enumi}.}
\setcounter{enumi}{12}
\tightlist
\item
  \begin{enumerate}
  \def\labelenumii{(\alph{enumii})}
  \tightlist
  \item
    A multiplication is defined on \(\mathbf{R}^3\) as follows:
  \end{enumerate}
\end{enumerate}

\[
(x, y, z). (x ^ {\prime}, y ^ {\prime}, z ^ {\prime}) = (x + x ^ {\prime} \cos z - y ^ {\prime} \sin z, y + x ^ {\prime} \sin z + y ^ {\prime} \cos z, z + z ^ {\prime}).
\]

Show that a solvable real Lie group \(\mathbf{G}\) is thus obtained such that \(\mathrm{DG} = \mathbf{R}^2\times \{0\}\) . The centre \(\mathbf{Z}\) of \(\mathbf{G}\) is \(\{0\} \times \{0\} \times 2\pi \mathbf{Z}\) .

\begin{enumerate}
\def\labelenumi{(\alph{enumi})}
\setcounter{enumi}{1}
\tightlist
\item
  For \((x,y,z)\in G\) , let \(\pi (x,y,z)\) be the mapping
\end{enumerate}

\[
(\lambda , \mu) \mapsto (\lambda \cos z - \mu \sin z + x, \lambda \sin z + \mu \cos z + y)
\]

of \(\mathbf{R}^2\) into \(\mathbf{R}^2\) . Show that \(\pi\) is a morphism of \(G\) onto a Lie subgroup \(\mathbf{G}'\) of the affine group of \(\mathbf{R}^2\) , generated by the translations and rotations of \(\mathbf{R}^2\) . Show that \(\operatorname{Ker} \pi = \mathbf{Z}\) .

\begin{enumerate}
\def\labelenumi{(\alph{enumi})}
\setcounter{enumi}{2}
\item
  We canonically identify \(\mathrm{L}(\mathrm{G}) = \mathrm{T}_{(0,0,0)}\mathrm{G}\) with \(\mathbf{R}^3\) . Let \((e_1, e_2, e_3)\) be the canonical basis of \(\mathbf{R}^3\) . Show that \([e_1, e_2] = 0\) , \([e_3, e_1] = e_2\) , \([e_3, e_2] = -e_1\) .
\item
  Show that, for \(c \neq 0\) ,
\end{enumerate}

\[
\exp_ {\mathrm{G}} (a, b, c) = \left(\frac {1}{c} (a \sin c + b \cos c - b), \frac {1}{c} (- a \cos c + b \sin c + a), c\right).
\]

Deduce that, for all \(u \in L(G)\) , \(Z \subset \exp(\mathbf{R}u)\) . Show that \(\exp_G: L(G) \to G\) is neither injective nor surjective.

\begin{enumerate}
\def\labelenumi{\arabic{enumi}.}
\setcounter{enumi}{13}
\tightlist
\item
  \begin{enumerate}
  \def\labelenumii{(\alph{enumii})}
  \tightlist
  \item
    Let \(\mathbf{G} = \mathbf{GL}(n, \mathbf{C})\) . Show that \(\exp_{\mathbf{G}}\) is surjective. (Use the holomorphic functional calculus of Spectral Theories, Chapter I, § 4, no. 8.)
  \end{enumerate}
\end{enumerate}

\begin{enumerate}
\def\labelenumi{(\alph{enumi})}
\setcounter{enumi}{1}
\tightlist
\item
  Let \(\mathbf{G}' = \mathbf{SL}(2, \mathbf{C})\) . Show that, if \(\sigma \in \mathbf{C}^*\) , the element \(\left( \begin{array}{cc} -1 & \sigma \\ 0 & -1 \end{array} \right)\)
\end{enumerate}

\(G'\) does not belong to the image of \(\exp_{G'}\) .

\begin{enumerate}
\def\labelenumi{\arabic{enumi}.}
\setcounter{enumi}{14}
\tightlist
\item
  \begin{enumerate}
  \def\labelenumii{(\alph{enumii})}
  \tightlist
  \item
    Let \(x \in \mathfrak{sl}(2, \mathbf{R})\) and \(\Delta = \det x\) . Show that
  \end{enumerate}
\end{enumerate}

\[
e ^ {x} = \operatorname{ch} \sqrt {- \Delta}. I + \frac {\operatorname{sh} \sqrt {- \Delta}}{\sqrt {- \Delta}} \cdot x \quad \mathrm{if} \Delta <   0
\]

\[
e ^ {x} = \cos {\sqrt {\Delta}}. I + \frac {\sin {\sqrt {\Delta}}}{\sqrt {\Delta}}. x \qquad \mathrm{if} \Delta > 0.
\]

\begin{enumerate}
\def\labelenumi{(\alph{enumi})}
\setcounter{enumi}{1}
\tightlist
\item
  Let \(\mathbf{H} = \mathbf{SL}(2, \mathbf{R})\) . Show that \(g = \begin{pmatrix} \lambda & 0 \\ 0 & \lambda^{-1} \end{pmatrix} \in \mathbf{H}\) is the image of \(\exp_{\mathbf{H}}\) if and only if \(\lambda > 0\) or \(\lambda = -1\) . If \(\lambda = 0\) , all the one-parameter subgroups containing \(g\) are equal.
\end{enumerate}

\begin{enumerate}
\def\labelenumi{\arabic{enumi}.}
\setcounter{enumi}{15}
\tightlist
\item
  Let G be a finite-dimensional Lie group. Show that there exists a basis \((x_{1},\ldots,x_{n})\) of L(G) such that, for all i, \(\exp(\mathbf{R}x_{i})\) is a Lie subgroup of G. (Use Spectral Theories, Chapter II, §2, Lemma 1.)
\end{enumerate}

¶ 17. (a) Let c denote a real solvable Lie algebra with a basis \((a, b, c)\) such that \([a, b] = c\) , \([a, c] = -b\) , \([b, c] = 0\) . Let d denote a real solvable Lie algebra with a basis \((a, b, c, d)\) such that \([a, b] = c\) , \([a, c] = -b\) , \([b, c] = d\) , \([a, d] = [b, d] = [c, d] = 0\) . Let g be a real solvable Lie algebra. If g contains non-zero elements x, y, z such that \([x, y] = z\) and \([x, z] = -y\) , then g contains a subalgebra isomorphic to c or d.~(Write \(a_{1} = y\) , \(b_{1} = z\) , \(c_{1} = [y, z]\) ; define \(a_{i}\) , \(b_{i}\) , \(c_{i}\) inductively by \(a_{i} = [a_{i-1}, c_{i-1}]\) , \(b_{i} = [b_{i-1}, c_{i-1}]\) , \(c_{i} = [a_{i}, b_{i}]\) . Consider the smallest integer k such that \(c_{k} = 0\) .)

\begin{enumerate}
\def\labelenumi{(\alph{enumi})}
\setcounter{enumi}{1}
\item
  Let g and \(g^{*}\) be real solvable Lie algebras and \(\phi\) a homomorphism of g onto \(g^{*}\) . If \(g^{*}\) contains a subalgebra isomorphic to c or d, g has the same property.
\item
  Let \(\mathfrak{h}\) be a complex solvable Lie algebra. Consider a Jordan-Hölder series for the adjoint representation of \(\mathfrak{h}\) ; the quotients of this series define 1-dimensional representations of \(\mathfrak{h}\) and hence linear forms on \(\mathfrak{h}\) . These linear forms, which depend only on \(\mathfrak{h}\) , are called the roots of \(\mathfrak{h}\) . If \(\mathfrak{h}'\) is a real solvable Lie algebra, the roots of \(\mathfrak{h}'\) are the restrictions to \(\mathfrak{h}'\) of the roots of \(\mathfrak{h}' \otimes_{\mathbf{R}} \mathbf{C}\) .
\end{enumerate}

Then let G be a finite-dimensional simply connected solvable real Lie group and g its Lie algebra. Show that the following conditions are equivalent: (α) exp₄ is injective; (β) exp₅ is surjective; (γ) exp₆ is bijective; (δ) exp₇ is an isomorphism of the analytic manifold L(G) onto the analytic manifold G; (ε) L(G) contains no subalgebra isomorphic to c or d; (ζ) there exists no quotient algebra of L(G) admitting a subalgebra isomorphic to c; (η) every root of g is of the form \(\phi + i\delta'\) where \(\phi\) , \(\phi'\) are in \(g^*\) and \(\phi'\) is proportional to \(\phi\) ; (θ) for all \(x \in g\) , the only pure imaginary eigenvalue (in C) of ad x is 0.†
